\chapter{Syntax: levels, expressions, declarations, environments}
\label{chap:syntax}

This chapter covers the object language of the whole verification. Everything that
later chapters state -- the typing judgement, the inductive specification, the
projection development, the kernel correspondence -- is stated about the universe
levels, expressions, declarations and environments defined here. The group contains
no judgements and no proofs about typing; it is the syntax plus its substitution
calculus, plus two self-contained extras (the translation from real Lean syntax, and
a complexity result about level equivalence).

The files are
\srcloc{Lean4Lean/Theory.lean}{1} (the aggregator module),
\srcloc{Lean4Lean/Theory/VLevel.lean}{1} (188 lines, universe levels),
\srcloc{Lean4Lean/Theory/VExpr.lean}{1} (1334 lines, expressions and substitution),
\srcloc{Lean4Lean/Theory/VDecl.lean}{1} (67 lines, declaration payloads),
\srcloc{Lean4Lean/Theory/VEnv.lean}{1} (57 lines, environments),
\srcloc{Lean4Lean/Theory/Quot.lean}{1} (21 lines, the quotient block),
\srcloc{Lean4Lean/Theory/Meta.lean}{1} (112 lines, elaboration glue) and
\srcloc{Lean4Lean/Theory/LevelSat.lean}{1} (450 lines, a coNP-hardness proof).
Of the 2235 lines, 1786 are master, 254 come from the \texttt{iota} branch and 195
from the \texttt{trproj} branch. The contributed lines are concentrated in three
places: a 372-line mixed block at the end of \texttt{VExpr.lean} (telescope and
spine readers, plus a substitution-based weakening section), more than half of
\texttt{VDecl.lean} (recursors and their computation rules), and the \texttt{pats}
field of \texttt{VEnv} with its extension operation.

Against Carneiro's thesis, \texttt{VLevel} and \texttt{VExpr} realise the grammar of
\texttt{axioms.tex} \S2.1 together with the constants of \S2.5, with two deliberate
deviations: universe variables are de Bruijn-style \texttt{Nat} indices rather than
names, and level order and equivalence are defined \emph{semantically} (pointwise over
all assignments) rather than by the syntactic derivation system of \S2.2. There is no
\texttt{let} and no projection in the syntax; both are eliminated at translation time
by \texttt{Theory/Meta.lean}, which is exactly why the \texttt{trproj} contribution has
to model kernel projections as a reduction relation elsewhere rather than as syntax.
\texttt{VDecl}/\texttt{VInductDecl} realise \S2.6.1, and the contributed
\texttt{VRecursor}/\texttt{VRecRule} realise \S2.6.3 and the constructor half of \S2.6.4.
(Section numbers throughout this chapter are the ones the thesis sources produce:
\texttt{intro.tex} is \S1 and \texttt{axioms.tex} is \S2, so \S2.5 is ``Definitions
($\delta$ reduction)'' and \S2.6 is ``Inductive types''.)
\texttt{LevelSat.lean} has no thesis counterpart at all; it is an original result of the
repository.

Two honest gaps against the thesis are visible already at this layer.
\texttt{VRecursor.k} carries the K-like-reduction flag but the thesis's second $\iota$
rule is not modelled; this one is recorded in the code's own docstrings, both on the
field (\srcloc{Lean4Lean/Theory/VDecl.lean}{32}) and on \texttt{VEnv.addRecRule}
(\srcloc{Lean4Lean/Theory/Inductive.lean}{251}). The second gap is not remarked on
anywhere: the quotient computation rule is still registered as a \texttt{VDefEq} rather
than through the new pattern registry, so the model carries two mechanisms for
computation rules, and \texttt{Theory/Quot.lean} carries no docstring at all.

Overall assessment: nothing in this group depends on a \texttt{sorry} -- every
declaration here, contributed or not, is a definition or a fully proved theorem, and no
declaration of these eight files has \texttt{sorryAx} in its axiom footprint -- and
documentation quality in the contributed regions is well above the surrounding master
code, with several docstrings explaining a design decision rather than restating a type.
The specification choices are good: total readers rather than partial ones, decidable
biconditional reformulations where decidability is wanted, and $n$-indexed general
forms where the projection development needs them. The main criticism is hygiene.
Seven contributed declarations in this group are dead -- five that nothing mentions at
all, and two that survive only by feeding one of those five -- and one \texttt{trproj}
refactor relocated an upstream lemma and left a six-declaration master block dead
without removing it, on a file that is otherwise shared verbatim with upstream.

\section{Module structure}

\begin{definition}[The \texttt{Lean4Lean.Theory} aggregator module]
\label{mod:theory-root}
\leanok
\contrib{mixed} \srcloc{Lean4Lean/Theory.lean}{1}
An import-only module, six lines long, that pulls in the whole idealised theory, in
file order: the environment lemmas, the strong typing judgement, unique typing,
Church-Rosser, the inductive-parameter development
\texttt{Theory.Typing.InductiveParams}, and head reduction. The fifth import (line 5,
\texttt{InductiveParams}) is the single line the \texttt{iota} branch added to this
file; \texttt{HeadReduction} on line 6 is master and is the last. A module is not a
declaration, so this node names no Lean constant.
\end{definition}

\section{Universe levels}

\begin{definition}[Universe level expressions \texttt{VLevel}]
\label{ind:vlevel}
\lean{Lean4Lean.VLevel}
\leanok
\srcloc{Lean4Lean/Theory/VLevel.lean}{8}
The grammar of universe levels: $0$, successor, $\max$, $\mathrm{imax}$, and
parameters \texttt{param i} indexed by a natural number. Universe variables are
represented by their index into the declaration's parameter list rather than by
name, so the thesis's substitution $\alpha[\bar\ell/\bar u]$ becomes list indexing.

\thesisref{axioms.tex \S2.1, the grammar $\ell ::= u \mid 0 \mid S\ell \mid \max(\ell,\ell) \mid \mathrm{imax}(\ell,\ell)$}
\end{definition}

\begin{definition}[Level well-formedness in a parameter count]
\label{def:vlevel-wf}
\lean{Lean4Lean.VLevel.WF,Lean4Lean.VLevel.decidable_WF}
\leanok
\uses{ind:vlevel}
\srcloc{Lean4Lean/Theory/VLevel.lean}{20}
\texttt{l.WF n} holds when every \texttt{param i} occurring in $l$ satisfies $i < n$,
i.e. when $l$ mentions only the first $n$ universe parameters. It is decidable, by the
obvious structural instance.

\thesisref{axioms.tex \S2.5, ``the universe variables in $\alpha$ and $\ell$ are contained in $\bar u$''}
\end{definition}

\begin{definition}[Evaluation of a level at a parameter assignment]
\label{def:vlevel-eval}
\lean{Lean4Lean.VLevel.eval}
\leanok
\uses{ind:vlevel}
\srcloc{Lean4Lean/Theory/VLevel.lean}{34}
\texttt{l.eval ls : Nat} interprets a level at an assignment \texttt{ls : List Nat} of
naturals to parameters, out-of-range parameters defaulting to $0$, with
$\mathrm{imax}$ interpreted by \texttt{Lean.Nat.imax}. This is the semantics against
which order and equivalence of levels are defined, rather than a syntactic proof
system; it is a deliberate simplification over the kernel's \texttt{Level.isEquiv},
and the reason a complexity statement about the latter can be made at all
(\ref{thm:levelsat-equiv-iff-unsat}).
\end{definition}

\begin{definition}[Order and equivalence of levels]
\label{def:vlevel-le-equiv}
\lean{Lean4Lean.VLevel.LE,Lean4Lean.VLevel.Equiv,Lean4Lean.VLevel.le_antisymm_iff,Lean4Lean.VLevel.equiv_def}
\leanok
\uses{def:vlevel-eval}
\srcloc{Lean4Lean/Theory/VLevel.lean}{41}
$a \le b$ means $a.\mathrm{eval}\;ls \le b.\mathrm{eval}\;ls$ for every assignment
$ls$, and $a \approx b$ means the two evaluation functions are equal;
\texttt{equiv\_def} restates the latter pointwise. Antisymmetry,
$a \approx b \leftrightarrow a \le b \wedge b \le a$, is \texttt{funext\_iff}
composed with \texttt{Nat.le\_antisymm\_iff} and \texttt{forall\_and}.

\thesisref{axioms.tex \S2.2, the algorithmic judgements $\ell \le \ell' + n$ and $\ell \equiv \ell'$, and the sort congruence rule $\ell\equiv\ell' \Rightarrow \mathsf{U}_\ell \equiv \mathsf{U}_{\ell'}$ that consumes them (the constant congruence rule is \S2.5)}
\end{definition}

\begin{lemma}[Order and congruence algebra of levels]
\label{fam:vlevel-lattice}
\lean{Lean4Lean.VLevel.le_refl,Lean4Lean.VLevel.le_trans,Lean4Lean.VLevel.zero_le,Lean4Lean.VLevel.le_succ,Lean4Lean.VLevel.succ_le_succ,Lean4Lean.VLevel.le_max_left,Lean4Lean.VLevel.le_max_right,Lean4Lean.VLevel.succ_congr,Lean4Lean.VLevel.max_congr,Lean4Lean.VLevel.imax_congr,Lean4Lean.VLevel.max_comm,Lean4Lean.VLevel.max_self,Lean4Lean.VLevel.zero_imax,Lean4Lean.VLevel.imax_zero,Lean4Lean.VLevel.imax_self,Lean4Lean.VLevel.imax_eq_zero,Lean4Lean.VLevel.IsNeverZero,Lean4Lean.VLevel.IsNeverZero.imax_eq_max}
\leanok
\uses{def:vlevel-le-equiv}
\srcloc{Lean4Lean/Theory/VLevel.lean}{45}
The routine algebra of $\le$ and $\approx$: reflexivity and transitivity, $0$ least,
$a \le \mathrm{succ}\,a$, the two $\max$ injections, congruence of
$\mathrm{succ}/\max/\mathrm{imax}$ under $\approx$, and the $\mathrm{imax}$
simplifications $\mathrm{imax}(0,a)\approx a$, $\mathrm{imax}(a,0)\approx 0$,
$\mathrm{imax}(a,a)\approx a$, $\mathrm{imax}(a,b)\approx 0 \leftrightarrow b\approx 0$,
and $\mathrm{imax}(a,b)\approx\max(a,b)$ when $b$ is never zero.
\end{lemma}
\begin{proof}
\uses{def:vlevel-eval}
\leanok
The $\le$ lemmas are term-mode applications of the matching \texttt{Nat} lemma under
the $\forall ls$ (\texttt{Nat.le\_refl}, \texttt{Nat.le\_trans},
\texttt{Nat.zero\_le}, \texttt{Nat.le\_succ}, \texttt{Nat.le\_max\_left/right}); the
$\approx$ lemmas are \texttt{simp [equiv\_def, eval]} plus the corresponding
\texttt{Nat.max}/\texttt{Lean.Nat.imax} rewrite. No \texttt{omega} is used anywhere in
this file. \texttt{imax\_eq\_zero} additionally needs a classical contradiction step
(\texttt{Decidable.byContradiction}).
\end{proof}

\begin{definition}[Level substitution and the identity substitution]
\label{def:vlevel-inst}
\lean{Lean4Lean.VLevel.inst,Lean4Lean.VLevel.params,Lean4Lean.VLevel.inst_inst,Lean4Lean.VLevel.inst_id,Lean4Lean.VLevel.inst_map_id,Lean4Lean.VLevel.eval_inst,Lean4Lean.VLevel.WF.inst,Lean4Lean.VLevel.inst_congr,Lean4Lean.VLevel.params_wf}
\leanok
\uses{ind:vlevel,def:vlevel-eval,def:vlevel-wf}
\srcloc{Lean4Lean/Theory/VLevel.lean}{115}
\texttt{l.inst ls} substitutes the levels \texttt{ls} for the parameters of $l$,
out-of-range parameters becoming $0$, and
$\texttt{params}\;n = [\texttt{param}\;0,\dots,\texttt{param}\;(n-1)]$ is the identity
substitution. The accompanying lemmas give functoriality
(\texttt{inst\_inst}), the identity law \texttt{inst\_id} for levels well-formed in
$n$ parameters, the retraction \texttt{inst\_map\_id}, commutation with
\texttt{eval}, preservation of well-formedness, and congruence under $\approx$
(pointwise on the substituted list).

\thesisref{axioms.tex \S2.5, $\tau_{\bar\ell}(c)=\alpha[\bar\ell/\bar u]$}
\end{definition}

\begin{definition}[Translation of kernel levels into the model]
\label{def:vlevel-oflevel}
\lean{Lean4Lean.VLevel.ofLevel,Lean4Lean.VLevel.WF.of_ofLevel,Lean4Lean.VLevel.WF.of_mapM_ofLevel}
\leanok
\uses{ind:vlevel,def:vlevel-wf}
\srcloc{Lean4Lean/Theory/VLevel.lean}{168}
\texttt{ofLevel ls l} partially translates a real \texttt{Lean.Level} into a
\texttt{VLevel}, resolving a named parameter to its index in \texttt{ls} and failing
on metavariables and on names not in \texttt{ls}. Any result is well-formed in
\texttt{ls.length}, and likewise for a list translated by \texttt{mapM}. This is the
bridge the \texttt{Verify} layer uses, and the only place where level \emph{names}
occur.
\end{definition}

\section{Expressions and the substitution calculus}

\begin{definition}[Expressions of the idealised theory \texttt{VExpr}]
\label{ind:vexpr}
\lean{Lean4Lean.VExpr}
\leanok
\srcloc{Lean4Lean/Theory/VExpr.lean}{7}
De Bruijn syntax with six constructors: \texttt{bvar i}, \texttt{sort u},
\texttt{const c us}, \texttt{app}, \texttt{lam} and \texttt{forallE}. Binders carry
only their domain type: there are no binder names, no \texttt{let}, no projections and
no literals. The absent constructs are eliminated at translation time
(\ref{def:meta-ofexpr}), which is why projections have to be modelled as a separate
reduction relation rather than as syntax.

\thesisref{axioms.tex \S2.1, the grammar $e ::= x \mid \mathsf{U}_\ell \mid e\,e \mid \lambda x{:}e.\,e \mid \forall x{:}e.\,e$, extended with $c_{\bar u}$ from \S2.5}
\end{definition}

\begin{definition}[Weakening \texttt{liftN} and its algebra]
\label{def:liftn}
\lean{Lean4Lean.liftVar,Lean4Lean.VExpr.liftN,Lean4Lean.VExpr.lift,Lean4Lean.VExpr.liftN_zero,Lean4Lean.VExpr.liftN'_liftN',Lean4Lean.VExpr.liftN'_liftN_lo,Lean4Lean.VExpr.liftN'_liftN_hi,Lean4Lean.VExpr.liftN_liftN,Lean4Lean.VExpr.liftN_succ,Lean4Lean.VExpr.liftN'_comm,Lean4Lean.VExpr.lift_liftN'}
\leanok
\uses{ind:vexpr}
\srcloc{Lean4Lean/Theory/VExpr.lean}{17}
\texttt{liftN n e k} adds $n$ to every de Bruijn index of $e$ that is at least the
cut-off $k$, descending with $k+1$ under binders; \texttt{lift} is \texttt{liftN 1}.
The lemmas compose two lifts under the side condition
$k_1 \le k_2 \le n_1 + k_1$ (\texttt{liftN'\_liftN'}), specialise it to the low and
high arrangements, and commute lifts at different cut-offs (\texttt{liftN'\_comm}).

\thesisref{typesys.tex, Lemma \texttt{thm:weak} (Weakening) -- this is its syntactic substrate}
\end{definition}

\begin{definition}[Closed expressions]
\label{def:closedn}
\lean{Lean4Lean.VExpr.ClosedN,Lean4Lean.VExpr.Closed,Lean4Lean.VExpr.ClosedN.mono,Lean4Lean.VExpr.ClosedN.liftN_eq,Lean4Lean.VExpr.ClosedN.lift_eq,Lean4Lean.VExpr.ClosedN.liftN,Lean4Lean.VExpr.ClosedN.liftN_eq_rev}
\leanok
\uses{ind:vexpr,def:liftn}
\srcloc{Lean4Lean/Theory/VExpr.lean}{99}
\texttt{e.ClosedN k} holds when every free de Bruijn index of $e$ is below $k$, with
\texttt{Closed} the default $k = 0$. Closedness is monotone in $k$, is preserved by
lifting (with the bound raised by $n$), and an expression closed below $k$ is fixed by
any lift with cut-off at least $k$.
\end{definition}

\begin{definition}[Decidability of \texttt{ClosedN}]
\label{inst:dec-closedn}
\lean{Lean4Lean.VExpr.decClosedN}
\leanok
\contrib{iota} \srcloc{Lean4Lean/Theory/VExpr.lean}{108}
\uses{def:closedn}
An instance \texttt{decClosedN} giving \texttt{Decidable (e.ClosedN k)} for all $e$ and
$k$, written as a direct structural definition of the decision witness rather than by
tactics: the \texttt{bvar} case is \texttt{Nat} comparison, \texttt{sort} and
\texttt{const} are \texttt{isTrue trivial}, and the three binary cases are
\texttt{instDecidableAnd}. It is needed because registering an $\iota$ rule branches on
\texttt{if h : ru.rhs.Closed}, so closedness is decided inside a \emph{definition}, not
only inside a proof.

\thesisref{no thesis counterpart (decidability is an implementation concern)}
\reviewnote{The sibling head-shape decidability instances were moved out of this file to
\texttt{Lean4Lean/Tests/ShapeDecide.lean} (commit \texttt{ba118fd}, ``file the head-shape
decidability instances with the tests'') while this one stayed. That file's module
docstring does give the general rationale -- the shape predicates are used as
propositions in the theory and only the tests run them -- but says nothing about why
\texttt{decClosedN} is the exception, and the exception is the interesting part, since it
is the one closedness test a \emph{definition} needs.}
\end{definition}

\begin{definition}[Universe instantiation of expressions]
\label{def:instl}
\lean{Lean4Lean.VExpr.instL,Lean4Lean.VExpr.LevelWF,Lean4Lean.VExpr.instL_instL,Lean4Lean.VExpr.instL_liftN,Lean4Lean.VExpr.LevelWF.instL_id,Lean4Lean.VExpr.levelWF_liftN,Lean4Lean.VExpr.ClosedN.instL}
\leanok
\uses{ind:vexpr,def:vlevel-inst}
\srcloc{Lean4Lean/Theory/VExpr.lean}{152}
\texttt{e.instL ls} applies the level substitution \texttt{ls} to every sort and to
every level argument of every constant in $e$, leaving the term structure untouched;
\texttt{e.LevelWF U} says every level occurring in $e$ is well-formed in $U$
parameters. Universe instantiation commutes with weakening, composes
(\texttt{instL\_instL}), preserves closedness, and is the identity on level-well-formed
terms at \texttt{VLevel.params U}.

\thesisref{axioms.tex \S2.5, the universe substitution applied to a constant's type and value}
\end{definition}

\begin{definition}[Single-variable instantiation and the lift/instantiation square]
\label{def:inst}
\lean{Lean4Lean.VExpr.instVar,Lean4Lean.VExpr.inst,Lean4Lean.VExpr.liftN_instN_lo,Lean4Lean.VExpr.liftN_instN_hi,Lean4Lean.VExpr.liftN_inst_hi,Lean4Lean.VExpr.lift_instN_lo,Lean4Lean.VExpr.lift_inst_hi,Lean4Lean.VExpr.inst_liftN,Lean4Lean.VExpr.inst_lift,Lean4Lean.VExpr.inst_liftN_lo,Lean4Lean.VExpr.inst_liftN_bvar}
\leanok
\uses{def:liftn,ind:vexpr}
\srcloc{Lean4Lean/Theory/VExpr.lean}{194}
\texttt{e.inst a k} substitutes $a$ for de Bruijn variable $k$ in $e$, shifting the
variables above $k$ down by one and lifting $a$ by $k$; the variable case
\texttt{instVar} is the three-way test $i < k$, $i = k$, $i > k$. The surrounding
lemmas are the standard commutation square between lifting and instantiation in the
\texttt{lo} arrangement (instantiation above the lift) and the \texttt{hi} arrangement
(instantiation below it), together with the cancellation
$(\texttt{liftN}\;1\;e\;k).\texttt{inst}\;a\;k = e$.

\thesisref{typesys.tex, Lemma \texttt{thm:subst} (Properties of substitution) -- its syntactic substrate}
\end{definition}

\begin{definition}[Telescope instantiation]
\label{def:insts}
\lean{Lean4Lean.VExpr.insts,Lean4Lean.VExpr.insts_append,Lean4Lean.VExpr.insts_sort,Lean4Lean.VExpr.insts_const,Lean4Lean.VExpr.insts_app,Lean4Lean.VExpr.insts_nil,Lean4Lean.VExpr.insts_cons}
\leanok
\uses{def:inst}
\srcloc{Lean4Lean/Theory/VExpr.lean}{333}
\texttt{e.insts as} substitutes a whole telescope by one \texttt{inst} at index $0$ per
element, the head of the list being substituted last, so that \texttt{insts (as ++ bs)}
is \texttt{insts as} after \texttt{insts bs}. The docstring argues for using a
substitution rather than an applied abstraction: statements built out of substituted
pieces decompose on the nose, since substitution commutes with application
syntactically, whereas an applied abstraction does so only up to definitional equality.

\thesisref{axioms.tex \S2.6.1, the telescope notation $f\;e : \beta[e/x]$}
\end{definition}

\begin{definition}[Inverting a lift: \texttt{unliftN} and \texttt{Skips}]
\label{def:unliftn-skips}
\lean{Lean4Lean.VExpr.unliftN,Lean4Lean.VExpr.Skips,Lean4Lean.VExpr.Skips',Lean4Lean.VExpr.skips_iff,Lean4Lean.VExpr.skips_iff_exists,Lean4Lean.VExpr.skips_add,Lean4Lean.VExpr.liftN_inj,Lean4Lean.VExpr.of_liftN_eq_liftN,Lean4Lean.VExpr.of_liftN_eq_instL,Lean4Lean.VExpr.Skips.of_liftN_hi,Lean4Lean.VExpr.Skips.of_instL}
\leanok
\uses{def:liftn,def:inst}
\srcloc{Lean4Lean/Theory/VExpr.lean}{367}
\texttt{unliftN e n k} undoes $n$ lifts at cut-off $k$ by instantiating
\texttt{default} $n$ times, and \texttt{e.Skips n k} is the equation
$\texttt{liftN}\;n\;(\texttt{unliftN}\;e\;n\;k)\;k = e$, i.e. ``$e$ is in the image of
\texttt{liftN n \_ k}''; \texttt{Skips'} is the equivalent structural characterisation.
The consequences are that \texttt{liftN} is injective, that two equal lifts factor
through a common preimage, and that skipping is inherited from a larger lift and from a
universe instantiation.
\end{definition}

\begin{lemma}[Commuting two instantiations; closedness under instantiation]
\label{fam:inst-inst}
\lean{Lean4Lean.VExpr.inst_instVar_hi,Lean4Lean.VExpr.inst_inst_hi,Lean4Lean.VExpr.inst0_inst_hi,Lean4Lean.VExpr.inst_instVar_lo,Lean4Lean.VExpr.inst_inst_lo,Lean4Lean.VExpr.instN_bvar0,Lean4Lean.VExpr.ClosedN.instN,Lean4Lean.VExpr.ClosedN.inst,Lean4Lean.VExpr.ClosedN.instN_eq,Lean4Lean.VExpr.instL_instN,Lean4Lean.VExpr.LevelWF.inst}
\leanok
\uses{def:inst,def:closedn}
\srcloc{Lean4Lean/Theory/VExpr.lean}{489}
The second half of the substitution calculus: two instantiations commute in both the
\texttt{hi} and \texttt{lo} cut-off arrangements; instantiation preserves
\texttt{ClosedN} (lowering the bound by one) and \texttt{LevelWF}; an expression closed
below $k$ is fixed by an instantiation at any index at least $k$; and universe
instantiation commutes with term instantiation.

\thesisref{typesys.tex, \texttt{item:subst\_ty} (Properties of substitution)}
\end{lemma}
\begin{proof}
\uses{def:instl,def:liftn}
\leanok
All by structural induction on the substituted term, the variable cases reducing to
arithmetic on \texttt{instVar} together with the lift lemmas of \ref{def:liftn}.
\end{proof}

\section{The \texttt{Lift} and \texttt{Subst} algebras}

\begin{definition}[The \texttt{Lift} renaming algebra]
\label{ind:lift}
\lean{Lean4Lean.Lift,Lean4Lean.Lift.skipN,Lean4Lean.Lift.consN,Lean4Lean.Lift.comp,Lean4Lean.Lift.dom,Lean4Lean.Lift.size,Lean4Lean.Lift.depth,Lean4Lean.Lift.inter,Lean4Lean.Lift.diff,Lean4Lean.Lift.Fixes,Lean4Lean.Lift.liftVar,Lean4Lean.Lift.liftVar_comp,Lean4Lean.Lift.Fixes.liftVar_eq,Lean4Lean.Lift.depth_succ}
\leanok
\uses{ind:vexpr}
\srcloc{Lean4Lean/Theory/VExpr.lean}{587}
A \texttt{Lift} is a finite description of an order-preserving injection of de Bruijn
indices, built from \texttt{refl}, \texttt{skip} (insert a binder) and \texttt{cons} (go
under a binder), acting on variables by \texttt{Lift.liftVar}. It carries a composition
monoid (\texttt{comp}, with $\texttt{liftVar}$ contravariantly functorial), a meet and a
difference (\texttt{inter}, \texttt{diff}), the measures \texttt{dom}, \texttt{size},
\texttt{depth}, and the predicate \texttt{Fixes k}, which says the renaming is the
identity below $k$. It generalises $\texttt{liftN}\;n\;\_\;k$, recovered as
$\texttt{consN}\;(\texttt{skipN}\;\texttt{refl}\;n)\;k$, to arbitrary interleavings.
\end{definition}

\begin{lemma}[\texttt{consN} fixes the binders it is pushed under]
\label{fam:lift-consn}
\lean{Lean4Lean.Lift.consN_fixes,Lean4Lean.Lift.consN_cons,Lean4Lean.Lift.liftVar_consN_lt,Lean4Lean.Lift.liftVar_consN_succ}
\leanok
\contrib{trproj} \srcloc{Lean4Lean/Theory/VExpr.lean}{744}
\uses{ind:lift}
Four facts about pushing a lift under $m$ binders: $(\rho.\texttt{consN}\;m)$ fixes $m$;
$(\texttt{cons}\;\rho).\texttt{consN}\;j = \rho.\texttt{consN}\;(j+1)$; the renaming is
the identity on indices below $m$; and it commutes with successor one binder at a time.

\thesisref{no thesis counterpart (de Bruijn bookkeeping)}
\end{lemma}
\begin{proof}
\uses{ind:lift}
\leanok
\texttt{consN\_fixes} and \texttt{consN\_cons} are structural recursions on $m$ and $j$;
\texttt{liftVar\_consN\_lt} is \texttt{consN\_fixes} fed to
\texttt{Lift.Fixes.liftVar\_eq}; \texttt{liftVar\_consN\_succ} is \texttt{rfl}.
\reviewnote{Of this four-lemma block only \texttt{consN\_cons} has a consumer (in
\texttt{piBinders\_lift'}, \ref{fam:pibinders-stability}).
\texttt{liftVar\_consN\_lt} and \texttt{liftVar\_consN\_succ} are unused anywhere in the
repository, which makes \texttt{consN\_fixes} dead in effect as well.}
\end{proof}

\begin{definition}[Weakening along a \texttt{Lift}]
\label{def:lift-prime}
\lean{Lean4Lean.VExpr.lift',Lean4Lean.VExpr.lift'_consN_skipN,Lean4Lean.VExpr.lift'_comp,Lean4Lean.VExpr.lift'_depth_zero,Lean4Lean.VExpr.lift'_refl,Lean4Lean.VExpr.lift_eq_lift',Lean4Lean.VExpr.instL_lift',Lean4Lean.VExpr.ClosedN.lift'_eq}
\leanok
\uses{ind:lift,def:liftn}
\srcloc{Lean4Lean/Theory/VExpr.lean}{763}
\texttt{e.lift' rho} renames the free variables of $e$ along $\rho$, descending with
\texttt{cons} under binders. It specialises to \texttt{liftN} at
$\texttt{consN}\;(\texttt{skipN}\;\texttt{refl}\;n)\;k$, is functorial for
\texttt{comp}, is the identity when $\rho$ has depth $0$, commutes with universe
instantiation, and fixes an expression closed below a level that $\rho$ fixes.
\end{definition}

\begin{definition}[General (simultaneous) substitutions]
\label{def:subst}
\lean{Lean4Lean.VExpr.Subst,Lean4Lean.VExpr.subst,Lean4Lean.VExpr.Subst.lift,Lean4Lean.VExpr.Subst.liftN,Lean4Lean.VExpr.Subst.lift_l,Lean4Lean.VExpr.Subst.lift_r,Lean4Lean.VExpr.subst_lift',Lean4Lean.VExpr.lift'_subst,Lean4Lean.VExpr.Subst.comp,Lean4Lean.VExpr.Subst.comp_lift,Lean4Lean.VExpr.subst_subst,Lean4Lean.VExpr.liftN_subst}
\leanok
\uses{def:lift-prime,ind:vexpr}
\srcloc{Lean4Lean/Theory/VExpr.lean}{794}
A \texttt{Subst} is a function \texttt{Nat} $\to$ \texttt{VExpr}; \texttt{e.subst sigma}
applies it, lifting under binders (\texttt{Subst.lift}, iterated as
\texttt{Subst.liftN}). \texttt{Subst.comp} composes two substitutions and
\texttt{subst\_subst} is its action law, $(e.\texttt{subst}\,\sigma).\texttt{subst}\,\sigma' = e.\texttt{subst}\,(\sigma.\texttt{comp}\,\sigma')$;
associativity and unitality of \texttt{comp} are not themselves stated, only
\texttt{comp\_lift}. A \texttt{Lift} is absorbed on either side through
\texttt{subst\_lift'} and \texttt{lift'\_subst}, which rewrite it into
\texttt{Subst.lift\_l} and \texttt{Subst.lift\_r}. This is the API the
\texttt{trproj} section at line 963 builds on.
\end{definition}

\begin{lemma}[Identity, cons-chains and the \texttt{inst}-as-\texttt{subst} bridge]
\label{fam:subst-basic}
\lean{Lean4Lean.VExpr.Subst.id,Lean4Lean.VExpr.Subst.head,Lean4Lean.VExpr.Subst.tail,Lean4Lean.VExpr.Subst.cons,Lean4Lean.VExpr.Subst.one,Lean4Lean.VExpr.subst_id,Lean4Lean.VExpr.id_lift,Lean4Lean.VExpr.Subst.lift_l_refl,Lean4Lean.VExpr.Subst.lift_l_skip,Lean4Lean.VExpr.lift_subst,Lean4Lean.VExpr.Subst.Fixes,Lean4Lean.VExpr.Subst.Fixes.lift,Lean4Lean.VExpr.ClosedN.subst_eq,Lean4Lean.VExpr.lift_subst_cons,Lean4Lean.VExpr.lift_subst_lift,Lean4Lean.VExpr.subst_inst,Lean4Lean.VExpr.inst_lift_cons,Lean4Lean.VExpr.instN_eq,Lean4Lean.VExpr.inst_eq}
\leanok
\uses{def:subst,def:inst}
\srcloc{Lean4Lean/Theory/VExpr.lean}{854}
The concrete substitutions \texttt{id}, \texttt{cons} and \texttt{one a} with their
head/tail algebra, the predicate \texttt{Subst.Fixes n} (identity below $n$, stable
under lifting), and the bridge
$\texttt{e.inst a k} = \texttt{e.subst ((Subst.one a).liftN k)}$ reducing
single-variable instantiation to a general substitution. A lift on the left of a
substitution walks off it one binder at a time, so a closing given as a \texttt{cons}
chain reaches \texttt{Subst.id} by \texttt{simp}.
\end{lemma}
\begin{proof}
\uses{def:lift-prime}
\leanok
\texttt{instN\_eq} is an induction on the term with a nested induction on the cut-off in
the variable case; the rest are \texttt{funext} plus a case split on the index, or short
inductions using \texttt{subst\_lift'}/\texttt{lift'\_subst}.
\end{proof}

\begin{definition}[Left inverse of a \texttt{Lift}, and injectivity of \texttt{lift'}]
\label{def:lift-inv}
\lean{Lean4Lean.Lift.inv,Lean4Lean.VExpr.lift_l_inv,Lean4Lean.VExpr.lift'_inj}
\leanok
\uses{def:subst,ind:lift}
\srcloc{Lean4Lean/Theory/VExpr.lean}{887}
\texttt{rho.inv} is the substitution that undoes $\rho$, the \texttt{skip} case
discarding the inserted variable to \texttt{default}; it satisfies
$\texttt{Subst.lift\_l}\;\rho\;\rho.\texttt{inv} = \texttt{Subst.id}$, whence
$\texttt{lift'}\;\_\;\rho$ is injective.
\end{definition}

\begin{lemma}[\texttt{Subst.Depth}, truncation and \texttt{lift\_r\_one}]
\label{fam:subst-trunc-orphans}
\lean{Lean4Lean.VExpr.Subst.Depth,Lean4Lean.VExpr.Subst.Depth.id,Lean4Lean.VExpr.Subst.Depth.one,Lean4Lean.VExpr.Subst.trunc,Lean4Lean.VExpr.Subst.lift_r_comm,Lean4Lean.VExpr.lift_r_one,Lean4Lean.VExpr.Subst.tail_eq_lift_l,Lean4Lean.VExpr.Subst.lift_r_tail}
\leanok
\uses{def:subst}
\srcloc{Lean4Lean/Theory/VExpr.lean}{796}
$\sigma.\texttt{Depth}\;n\;n'$ says $\sigma$ is the shift by $n-n'$ above $n'$;
\texttt{Subst.trunc} truncates a substitution to such a shift, and
\texttt{Subst.lift\_r\_comm} and \texttt{lift\_r\_one} push a \texttt{Lift} through a
substitution of known depth. On master this block existed in order to prove
\texttt{lift'\_inst\_hi}.
\end{lemma}
\begin{proof}
\uses{def:lift-prime}
\leanok
\texttt{lift\_r\_comm} is \texttt{funext} plus an induction computing
$(\rho.\texttt{consN}\;n).\texttt{liftVar}$; \texttt{lift\_r\_one} instantiates it at
\texttt{Subst.Depth.one} and simplifies the truncation away.
\reviewnote{The \texttt{trproj} branch re-derived \texttt{lift'\_inst\_hi} from
\ref{thm:lift-instn-hi} and left this block behind. \texttt{lift\_r\_one}
(\srcloc{Lean4Lean/Theory/VExpr.lean}{919}) now has no consumer at all, and the whole
chain that fed it died with it: \texttt{Subst.lift\_r\_comm} (911), \texttt{Subst.trunc}
(908), \texttt{Subst.Depth.one} (885), \texttt{Subst.Depth.id} (858) and
\texttt{Subst.Depth} itself (796) are now mentioned only inside this block -- six
declarations. \texttt{Subst.tail\_eq\_lift\_l} was already unused on master. Only
\texttt{Subst.lift\_r\_tail} is still used, from
\srcloc{Lean4Lean/Theory/Typing/Lemmas.lean}{1005}. All six survive verbatim in
\texttt{Experimental/SExpr.lean}, where the same API is proved out and
\texttt{lift\_r\_one} is still consumed.}
\end{proof}

\section{Weakening and instantiation as substitutions}

The \texttt{trproj} branch opened this section of \texttt{VExpr.lean} with a design
statement (\srcloc{Lean4Lean/Theory/VExpr.lean}{963}): since \texttt{lift'} and
\texttt{inst} are both instances of \texttt{subst}, a lemma about a term builder need
only be proved once against \texttt{subst}, and the weakening and instantiation forms
are read off by rewriting, \texttt{Lift.consN} matching \texttt{Subst.liftN} binder for
binder. The rest of the section is that programme carried out.

\begin{theorem}[Weakening is a substitution]
\label{thm:lift-eq-subst}
\lean{Lean4Lean.VExpr.lift'_eq_subst}
\leanok
\contrib{trproj} \srcloc{Lean4Lean/Theory/VExpr.lean}{970}
\uses{def:subst,def:lift-prime}
$\texttt{e.lift'}\;\rho = \texttt{e.subst}\;(\texttt{Subst.lift\_l}\;\rho\;\texttt{Subst.id})$:
weakening along $\rho$ is the substitution that renames each variable along $\rho$.

\thesisref{no thesis counterpart (an internal refactoring lemma)}
\end{theorem}
\begin{proof}
\uses{fam:subst-basic,def:subst}
\leanok
\texttt{subst\_id.symm.trans subst\_lift'}: run the identity substitution on
$\texttt{e.lift'}\;\rho$, then absorb the lift on the left by \texttt{subst\_lift'}.
This is the keystone of the section and is used from \texttt{Theory/Proj.lean} as well
as inside this file.
\end{proof}

\begin{lemma}[\texttt{Subst.liftN} bookkeeping]
\label{fam:subst-liftn-algebra}
\lean{Lean4Lean.VExpr.Subst.lift_liftN,Lean4Lean.VExpr.Subst.liftN_liftN,Lean4Lean.VExpr.Subst.lift_lift_l_id,Lean4Lean.VExpr.Subst.liftN_lift_l_id,Lean4Lean.VExpr.Subst.Fixes.liftN}
\leanok
\contrib{trproj} \srcloc{Lean4Lean/Theory/VExpr.lean}{973}
\uses{def:subst,def:lift-prime}
Five unfolding lemmas for iterated lifting of a substitution:
$\sigma.\texttt{lift}.\texttt{liftN}\;n = \sigma.\texttt{liftN}\;(n+1)$;
$(\sigma.\texttt{liftN}\;k).\texttt{liftN}\;j = \sigma.\texttt{liftN}\;(k+j)$;
$\texttt{Subst.lift\_l}\;\rho\;\texttt{id}$ commutes with \texttt{lift} and
\texttt{liftN}, turning $\rho$ into $\rho.\texttt{cons}$ and $\rho.\texttt{consN}\;j$;
and $(\sigma.\texttt{liftN}\;n)$ fixes $n$.

\thesisref{no thesis counterpart (de Bruijn bookkeeping)}
\end{lemma}
\begin{proof}
\uses{fam:subst-basic}
\leanok
\texttt{Subst.lift\_liftN} and \texttt{Subst.liftN\_liftN} are structural recursions on
the number of lifts whose step is \texttt{congrArg Subst.lift} on the induction
hypothesis; \texttt{Subst.Fixes.liftN} recurses through \texttt{Subst.Fixes.lift};
\texttt{Subst.liftN\_lift\_l\_id} recurses by rewriting with
\texttt{Subst.lift\_lift\_l\_id}, which is itself not a recursion at all but the two
rewrites \texttt{Subst.lift\_l\_lift} and \texttt{id\_lift}.
\texttt{Subst.liftN\_lift\_l\_id} is the one that makes \texttt{Lift.consN} and
\texttt{Subst.liftN} line up binder for binder, which is what the section docstring
promises; all five are consumed from \texttt{Theory/Proj.lean} (lines 171, 177, 189,
343 and 368).
\end{proof}

\begin{lemma}[Substitution past a one-variable instantiation, at the level of substitutions]
\label{thm:subst-comp-liftn-one}
\lean{Lean4Lean.VExpr.Subst.comp_liftN_one}
\leanok
\contrib{trproj} \srcloc{Lean4Lean/Theory/VExpr.lean}{993}
\uses{def:subst,fam:subst-basic}
For every $n$,
$\mathrm{comp}\,((\mathrm{one}\,P).\mathrm{liftN}\,n)\,(\sigma.\mathrm{liftN}\,n) = \mathrm{comp}\,(\sigma.\mathrm{liftN}\,(n{+}1))\,((\mathrm{one}\,(P.\mathrm{subst}\,\sigma)).\mathrm{liftN}\,n)$,
i.e. the substitution lemma stated between substitutions rather than between terms.

\thesisref{no thesis counterpart}
\end{lemma}
\begin{proof}
\uses{fam:subst-basic,def:subst}
\leanok
Induction on $n$. The base case is \texttt{funext} followed by a case split on the
variable, the successor branch closing with \texttt{lift\_subst\_cons} and
\texttt{subst\_id}; the step rewrites by \texttt{Subst.comp\_lift} on both sides and
applies the induction hypothesis. Stating it at the \texttt{Subst} level is the right
call: all the arithmetic lives here, and \ref{thm:subst-instn} below is then a rewrite.
Its only consumer is that lemma.
\end{proof}

\begin{theorem}[Substitution commutes with instantiation at index $n$]
\label{thm:subst-instn}
\lean{Lean4Lean.VExpr.subst_instN}
\leanok
\contrib{trproj} \srcloc{Lean4Lean/Theory/VExpr.lean}{1009}
\uses{def:subst,def:inst}
$(e.\mathrm{inst}\,a\,n).\mathrm{subst}\,(\sigma.\mathrm{liftN}\,n) = (e.\mathrm{subst}\,(\sigma.\mathrm{liftN}\,(n{+}1))).\mathrm{inst}\,(a.\mathrm{subst}\,\sigma)\,n$,
for all $e$, $a$, $\sigma$ and $n$.

\thesisref{typesys.tex, Lemma \texttt{thm:subst} in its general ($n$-indexed) form}
\end{theorem}
\begin{proof}
\uses{thm:subst-comp-liftn-one,fam:subst-basic}
\leanok
Rewrite both instantiations to substitutions by \texttt{instN\_eq}, collapse the two
composites by \texttt{subst\_subst}, and close with \ref{thm:subst-comp-liftn-one}.
\reviewnote{Master already proves \texttt{subst\_inst}
(\srcloc{Lean4Lean/Theory/VExpr.lean}{952}), which is exactly the case $n = 0$, and
keeps its independent hand proof; the contribution did not fold it in. The docstring
says ``under \texttt{m} binders'' while the statement binds $n$.}
\end{proof}

\begin{theorem}[Weakening under $m$ binders commutes with instantiation at index $m$]
\label{thm:lift-instn-hi}
\lean{Lean4Lean.VExpr.lift'_instN_hi,Lean4Lean.VExpr.lift'_inst_hi}
\leanok
\contrib{trproj} \srcloc{Lean4Lean/Theory/VExpr.lean}{1014}
\uses{def:lift-prime,def:inst,ind:lift}
$(e_1.\mathrm{inst}\,e_2\,m).\mathrm{lift}'\,(\rho.\mathrm{consN}\,m) = (e_1.\mathrm{lift}'\,(\rho.\mathrm{consN}\,(m{+}1))).\mathrm{inst}\,(e_2.\mathrm{lift}'\,\rho)\,m$,
with \texttt{lift'\_inst\_hi} the case $m = 0$.

\thesisref{typesys.tex, the weakening/substitution interaction of Lemma \texttt{thm:subst}}
\end{theorem}
\begin{proof}
\uses{thm:lift-eq-subst,fam:subst-liftn-algebra,thm:subst-instn}
\leanok
One \texttt{simp only} call: read both sides as substitutions by
\ref{thm:lift-eq-subst}, rewrite $\texttt{Subst.lift\_l}\;\rho\;\texttt{id}$ lifted $m$
times backwards through \texttt{Subst.liftN\_lift\_l\_id}, and apply
\ref{thm:subst-instn}.
\reviewnote{\texttt{lift'\_inst\_hi} already existed on master at line 903 with a
one-line proof from \texttt{lift\_r\_one}. Moving it 115 lines down and re-deriving it
is arguably cleaner, but it creates merge-conflict surface on a file otherwise shared
verbatim with upstream and orphans the block at \ref{fam:subst-trunc-orphans}. The
general form \texttt{lift'\_instN\_hi} has no consumer outside this pair.}
\end{proof}

\begin{lemma}[Inserting a binder at depth $i$ commutes with a shifted substitution]
\label{thm:liftn-subst-liftn}
\lean{Lean4Lean.VExpr.liftN_subst_liftN,Lean4Lean.VExpr.substVar_liftN_lift}
\leanok
\contrib{trproj} \srcloc{Lean4Lean/Theory/VExpr.lean}{1032}
\uses{def:subst,def:liftn}
$(e.\mathrm{liftN}\,1\,i).\mathrm{subst}\,(\sigma.\mathrm{liftN}\,(i{+}1)) = (e.\mathrm{subst}\,(\sigma.\mathrm{liftN}\,i)).\mathrm{liftN}\,1\,i$, with
\texttt{substVar\_liftN\_lift} the variable case.

\thesisref{typesys.tex, Lemma \texttt{thm:subst}, generalising master's \texttt{lift\_subst\_lift}}
\end{lemma}
\begin{proof}
\uses{def:liftn}
\leanok
Structural recursion on $e$, the binder cases raising $i$ by one; the \texttt{bvar} case
is \texttt{substVar\_liftN\_lift}, itself a recursion on the pair of the depth $i$ and
the index, using \texttt{liftVar\_succ} and \texttt{lift\_liftN'}.
\reviewnote{The docstring correctly identifies the case $i = 0$ as master's
\texttt{lift\_subst\_lift} (\srcloc{Lean4Lean/Theory/VExpr.lean}{947}) but does not
re-derive it, so that specialisation keeps a separate hand proof; the lemma itself has
one consumer, in \texttt{Theory/Proj.lean}.}
\end{proof}

\section{Telescopes and application spines}

The remainder of \texttt{VExpr.lean} is the \texttt{iota} branch's toolkit for talking
about the \emph{shape} of a recursor type, a constructor type and an $\iota$ redex,
together with \texttt{trproj} stability lemmas saying those shapes survive substitution.
All the readers are total functions rather than \texttt{Option}-valued ones, which is
what lets the well-formedness predicates built on them be decidable and hence testable
against the real kernel.

\begin{definition}[Application spines \texttt{mkApps}]
\label{def:mkapps}
\lean{Lean4Lean.VExpr.mkApps,Lean4Lean.VExpr.mkApps_nil,Lean4Lean.VExpr.mkApps_cons,Lean4Lean.VExpr.mkApps_append}
\leanok
\contrib{iota} \srcloc{Lean4Lean/Theory/VExpr.lean}{1051}
\uses{ind:vexpr}
$\texttt{f.mkApps}\;[a_0,\dots,a_n] = f\;a_0\dots a_n$, the left fold of \texttt{app},
with the \texttt{nil} and \texttt{cons} simp lemmas and associativity over list append.
It was introduced so that recursor redexes and constructor applications can be written
uniformly; \texttt{Theory/Inductive.lean} and \texttt{Theory/Proj.lean} are built on it.

\thesisref{axioms.tex \S2.6.1, the telescope application $f\;e_1\dots e_n$}
\end{definition}

\begin{lemma}[Spines commute with substitution, weakening and instantiation]
\label{fam:mkapps-stability}
\lean{Lean4Lean.VExpr.mkApps_subst,Lean4Lean.VExpr.mkApps_lift',Lean4Lean.VExpr.mkApps_inst,Lean4Lean.VExpr.mkApps_instL}
\leanok
\contrib{trproj} \srcloc{Lean4Lean/Theory/VExpr.lean}{1059}
\uses{def:mkapps,def:subst,def:inst,def:instl,def:lift-prime}
Each of \texttt{subst}, \texttt{lift'}, \texttt{inst} and \texttt{instL} distributes
over \texttt{mkApps}, applying the operation to the head and mapping it over the
argument list.

\thesisref{no thesis counterpart (structural congruences)}
\end{lemma}
\begin{proof}
\uses{thm:lift-eq-subst,fam:subst-basic}
\leanok
\texttt{mkApps\_subst} is one induction on the argument list; \texttt{mkApps\_lift'} and
\texttt{mkApps\_inst} are then \texttt{simp only} corollaries via
\ref{thm:lift-eq-subst} and \texttt{instN\_eq}. \texttt{mkApps\_instL} needs its own
induction because \texttt{instL} is not a \texttt{Subst}. This is the section's design
statement paying off: one induction, three corollaries, and all four are consumed --
\texttt{subst}, \texttt{inst} and \texttt{instL} from \texttt{Theory/Proj.lean} (lines
300, 309, 310) and the \texttt{Verify} layer, \texttt{mkApps\_lift'} from
\srcloc{Lean4Lean/Verify/Typing/Lemmas.lean}{656}.
\end{proof}

\begin{definition}[A descending run of de Bruijn variables]
\label{def:bvarsdesc}
\lean{Lean4Lean.VExpr.bvarsDesc,Lean4Lean.VExpr.bvarsDesc_length,Lean4Lean.VExpr.getElem_bvarsDesc}
\leanok
\contrib{mixed} \srcloc{Lean4Lean/Theory/VExpr.lean}{1081}
\uses{ind:vexpr}
$\texttt{bvarsDesc}\;lo\;n = [\texttt{bvar}\;(lo+n-1),\dots,\texttt{bvar}\;lo]$, the
argument list with which a constructor or recursor applies the $n$ variables of a
telescope it has just bound; its length is $n$ and its $t$-th element is
$\texttt{bvar}\;(lo + (n-1-t))$. The descending convention is the one forced by de
Bruijn telescopes and is documented. Definition and length lemma are \texttt{iota}
(lines 1081 and 1083); \texttt{getElem\_bvarsDesc} (line 1085) was added on
\texttt{trproj} for the projection development.

\thesisref{axioms.tex \S2.6.1, the variables $b$ of a constructor telescope}
\end{definition}

\begin{definition}[Reading the leading Pi-telescope]
\label{def:pi-telescope}
\lean{Lean4Lean.VExpr.piArity,Lean4Lean.VExpr.piBody,Lean4Lean.VExpr.piBinders,Lean4Lean.VExpr.piBinders_length}
\leanok
\contrib{iota} \srcloc{Lean4Lean/Theory/VExpr.lean}{1090}
\uses{ind:vexpr}
Three total functions on \texttt{VExpr}: \texttt{piArity} counts the leading
\texttt{forallE} binders, \texttt{piBody} strips them all, and \texttt{piBinders} lists
their domains outermost first, each still in its own binder context; the list has length
\texttt{piArity}. Totality rather than an \texttt{Option} is the right choice here,
because the shape predicates that consume these are conjunctions that pin the arity
separately.

\thesisref{axioms.tex \S2.6.1, the telescope notation $(x{::}\alpha)$ and $\forall x{::}\alpha.\,\beta$}
\end{definition}

\begin{lemma}[Pi-telescopes under weakening and universe instantiation]
\label{fam:pibinders-stability}
\lean{Lean4Lean.VExpr.piBinders_lift',Lean4Lean.VExpr.piBinders_instL,Lean4Lean.VExpr.piArity_instL}
\leanok
\contrib{trproj} \srcloc{Lean4Lean/Theory/VExpr.lean}{1108}
\uses{def:pi-telescope,def:lift-prime,def:instl}
$(T.\texttt{lift'}\;\rho).\texttt{piBinders}$ is $T.\texttt{piBinders}$ with the $j$-th
domain weakened along $\rho.\texttt{consN}\;j$; universe instantiation maps over the
binders unchanged and therefore preserves \texttt{piArity}.

\thesisref{no thesis counterpart (stability of a syntactic reader)}
\end{lemma}
\begin{proof}
\uses{fam:lift-consn}
\leanok
Structural recursion following the \texttt{forallE} spine, using
\texttt{List.mapIdx\_cons} and \texttt{Lift.consN\_cons} to shift the index; the
non-\texttt{forallE} cases are \texttt{rfl}. Note the asymmetry with
\ref{fam:ctorheaded-stability}: \texttt{lift'} and \texttt{instL} cannot create new
leading binders, whereas \texttt{inst} can, which is why the \texttt{inst} version needs
a \texttt{CtorHeaded} hypothesis.
\end{proof}

\begin{definition}[Reading the leading lambda-telescope]
\label{def:lam-telescope}
\lean{Lean4Lean.VExpr.lamArity,Lean4Lean.VExpr.lamBody,Lean4Lean.VExpr.lamBinders,Lean4Lean.VExpr.lamBinders_length,Lean4Lean.VExpr.foldr_lam_lamBinders}
\leanok
\contrib{mixed} \srcloc{Lean4Lean/Theory/VExpr.lean}{1123}
\uses{ind:vexpr}
The $\lambda$-analogues of the $\Pi$-readers: \texttt{lamArity}, \texttt{lamBody} and
\texttt{lamBinders}, with the length lemma and the reconstruction
$\texttt{e.lamBinders.foldr lam e.lamBody} = e$. \texttt{lamArity} and \texttt{lamBody}
are \texttt{iota} (lines 1123 and 1128); \texttt{lamBinders}, its length lemma and the
reconstruction are \texttt{trproj} (lines 1133 to 1143).

\thesisref{no thesis counterpart (a reader for $\iota$-rule reducts)}
\reviewnote{The consumer of all three \texttt{trproj} additions is a single proof in the
verification layer, \srcloc{Lean4Lean/Verify/Environment/Lemmas.lean}{1128}: it builds the
$\beta$-redex of an $\iota$ rule out of \texttt{rhs.lamBinders} and closes the
reconstruction with \texttt{foldr\_lam\_lamBinders} (lines 1133 and 1134).
\texttt{lamBinders\_length} is named nowhere, but it is \texttt{@[simp]} and the arity
side condition at line 1132 (\texttt{simp [hla]} against
\texttt{hla : rhs.lamArity = ...}) is exactly what it discharges, so it should not be
called dead. This is the one \texttt{lamBinders} consumer in the repository; nothing in
\texttt{Theory/Proj.lean} uses these three.}
\end{definition}

\begin{definition}[Decomposing an application spine]
\label{def:app-spine}
\lean{Lean4Lean.VExpr.getAppFn,Lean4Lean.VExpr.getAppArgs,Lean4Lean.VExpr.getAppFn_app,Lean4Lean.VExpr.getAppArgs_app,Lean4Lean.VExpr.getAppFn_mkApps,Lean4Lean.VExpr.getAppArgs_mkApps,Lean4Lean.VExpr.mkApps_getAppFn_getAppArgs}
\leanok
\contrib{iota} \srcloc{Lean4Lean/Theory/VExpr.lean}{1146}
\uses{def:mkapps}
\texttt{getAppFn} returns the head of an application spine and \texttt{getAppArgs} its
arguments, first argument first. They compute on \texttt{mkApps} (the head is the head of
the inner function, the arguments are appended) and invert it:
$\texttt{e.getAppFn.mkApps e.getAppArgs} = e$. The names deliberately mirror
\texttt{Lean.Expr.getAppFn} and \texttt{getAppArgs}, since the \texttt{Verify} layer has
to match these against the kernel's own spine functions.

\thesisref{axioms.tex \S2.6.1, reading a target $t\;e$ as a head applied to arguments}
\end{definition}

\begin{theorem}[Prefix characterisation of an application spine]
\label{thm:eq-mkapps-append}
\lean{Lean4Lean.VExpr.eq_mkApps_append_iff,Lean4Lean.VExpr.eq_mkApps_append_length_iff}
\leanok
\contrib{iota} \srcloc{Lean4Lean/Theory/VExpr.lean}{1174}
\uses{def:app-spine,def:mkapps}
$\exists\,\mathit{rest},\; e = f.\texttt{mkApps}\,(\mathit{pre} \mathbin{+\!\!+} \mathit{rest})$
holds if and only if $e$ and $f$ have the same spine head and
$f.\texttt{getAppArgs} \mathbin{+\!\!+} \mathit{pre}$ is a prefix of
$e.\texttt{getAppArgs}$. The second version additionally pins
$|\mathit{rest}| = n$ by the equation
$|e.\texttt{getAppArgs}| = |f.\texttt{getAppArgs}| + |\mathit{pre}| + n$.

\thesisref{no thesis counterpart (a decidability-oriented reformulation)}
\end{theorem}
\begin{proof}
\uses{def:app-spine}
\leanok
Forward by computing both spine projections of
$f.\texttt{mkApps}\,(\mathit{pre} \mathbin{+\!\!+} \mathit{rest})$; backward by rewriting
\texttt{mkApps\_getAppFn\_getAppArgs e} along the head equation and the prefix witness,
the length version then closing by \texttt{omega}. These are the workhorses that turn
the existential shape predicates of \texttt{Theory/Inductive.lean} into decidable ones:
the existential over a list of expressions on the left becomes, on the right, a spine
equality plus a \texttt{List.IsPrefix} and a length equation, all decidable once
\texttt{VExpr} has \texttt{DecidableEq}. The two instances built directly on them are at
\srcloc{Lean4Lean/Tests/ShapeDecide.lean}{20}.
\end{proof}

\begin{definition}[Head-constructor tests and their reflections]
\label{def:head-tests}
\lean{Lean4Lean.VExpr.isSort,Lean4Lean.VExpr.isBvar,Lean4Lean.VExpr.isConst,Lean4Lean.VExpr.isSort_iff,Lean4Lean.VExpr.isBvar_iff,Lean4Lean.VExpr.isConst_iff}
\leanok
\contrib{iota} \srcloc{Lean4Lean/Theory/VExpr.lean}{1200}
\uses{ind:vexpr}
Boolean tests for the \texttt{sort}, \texttt{bvar} and \texttt{const} head
constructors, each with an \texttt{iff} lemma reflecting the test into the corresponding
existential.

\thesisref{no thesis counterpart (test support)}
\reviewnote{The only consumers are the \texttt{Decidable} instances in
\texttt{Lean4Lean/Tests/ShapeDecide.lean}, so this is test-support machinery living in a
theory file. Also, the docstring ``Head-constructor tests, with their reflections into
existentials'' is attached to \texttt{isSort} alone while describing all six
declarations.}
\end{definition}

\begin{definition}[\texttt{RecHeaded}: the shape of a recursor type]
\label{def:recheaded}
\lean{Lean4Lean.VExpr.RecHeaded}
\leanok
\contrib{iota} \srcloc{Lean4Lean/Theory/VExpr.lean}{1219}
\uses{def:pi-telescope,def:app-spine}
$\texttt{ty.RecHeaded} := \exists k,\; \texttt{ty.piBody.getAppFn} = \texttt{bvar}\;k$:
the $\Pi$-body of \texttt{ty} is an application headed by a bound variable. That is the
shape of a recursor's type and of a minor premise, whose targets are applications of a
motive binder. It is the cheap, purely syntactic discriminator that separates recursor
constants from constructor constants.

\thesisref{axioms.tex \S2.6.3, the recursor type $\forall C{:}\kappa.\,\forall e{::}\varepsilon.\,\forall a{::}\alpha.\,\forall z{:}P\,a.\;C\,a\,z$, whose body is headed by the motive variable $C$}
\end{definition}

\begin{definition}[\texttt{CtorHeaded}, and its disjointness from \texttt{RecHeaded}]
\label{def:ctorheaded}
\lean{Lean4Lean.VExpr.CtorHeaded,Lean4Lean.VExpr.RecHeaded.not_ctorHeaded}
\leanok
\contrib{iota} \srcloc{Lean4Lean/Theory/VExpr.lean}{1223}
\uses{def:pi-telescope,def:app-spine,def:recheaded}
$\texttt{ty.CtorHeaded} := \exists I\,us,\; \texttt{ty.piBody.getAppFn} = \texttt{const}\;I\;us$:
the $\Pi$-body is an application headed by a \emph{constant}, the shape of a
constructor's type, whose target is a type-former application. Since a bound variable is
not a constant, no type is both \texttt{RecHeaded} and \texttt{CtorHeaded}; the proof
rewrites the \texttt{bvar} witness into the \texttt{const} equation and closes by
\texttt{cases}. That disjointness lemma is the load-bearing one: it is what rules out an
$\iota$ redex whose head is simultaneously a recursor and a constructor
(\ref{thm:indparams-patsiota-rec-ne-ctor}).

\thesisref{axioms.tex \S2.6.1, the constructor judgement $\Gamma;t{:}F\vdash \alpha\;\mathsf{ctor}$, whose base case is $t\;e$}
\end{definition}

\begin{lemma}[\texttt{CtorHeaded} passes to the body of a Pi]
\label{thm:ctorheaded-forallE}
\lean{Lean4Lean.VExpr.CtorHeaded.forallE}
\leanok
\contrib{trproj} \srcloc{Lean4Lean/Theory/VExpr.lean}{1232}
\uses{def:ctorheaded}
If $\texttt{forallE}\;A\;B$ is \texttt{CtorHeaded} then so is $B$.

\thesisref{no thesis counterpart}
\end{lemma}
\begin{proof}
\uses{def:pi-telescope}
\leanok
The proof term is \texttt{h}: the statement is definitionally trivial, because
$(\texttt{forallE}\;A\;B).\texttt{piBody}$ reduces to $B.\texttt{piBody}$.
\reviewnote{The lemma has no consumer anywhere in the repository, and its docstring
states a different proposition than the statement: ``a \texttt{CtorHeaded} type is not a
$\Pi$-telescope ending in a variable, so instantiation cannot create new leading
binders'' is the motivation for \texttt{piArity\_inst\_of\_ctorHeaded}, which lives some
forty lines below.}
\end{proof}

\begin{lemma}[\texttt{CtorHeaded} is stable under instantiation, and then so is the arity]
\label{fam:ctorheaded-stability}
\lean{Lean4Lean.VExpr.getAppFn_inst_const,Lean4Lean.VExpr.CtorHeaded.inst,Lean4Lean.VExpr.getAppFn_instL_const,Lean4Lean.VExpr.CtorHeaded.instL,Lean4Lean.VExpr.piBinders_inst_of_ctorHeaded,Lean4Lean.VExpr.piArity_inst_of_ctorHeaded}
\leanok
\contrib{trproj} \srcloc{Lean4Lean/Theory/VExpr.lean}{1234}
\uses{def:ctorheaded,def:pi-telescope,def:inst,def:instl}
A spine whose head is a constant keeps that head under \texttt{inst}, and keeps it up to
level substitution under \texttt{instL}; hence \texttt{CtorHeaded} is preserved by both.
Consequently instantiating a \texttt{CtorHeaded} type maps its $\Pi$-binders pointwise,
the $j$-th binder being instantiated at index $k+j$, and leaves \texttt{piArity}
unchanged.

\thesisref{no thesis counterpart (stability needed by the projection development)}
\end{lemma}
\begin{proof}
\uses{def:app-spine}
\leanok
All by structural recursion following the spine and then the telescope, the impossible
head shapes discharged by \texttt{nomatch};
\texttt{piArity\_inst\_of\_ctorHeaded} is \texttt{piBinders\_inst\_of\_ctorHeaded}
followed by \texttt{List.length\_mapIdx}.
\reviewnote{\texttt{getAppFn\_instL\_const}
(\srcloc{Lean4Lean/Theory/VExpr.lean}{1248}) and \texttt{CtorHeaded.instL}
(\srcloc{Lean4Lean/Theory/VExpr.lean}{1256}) are dead: nothing outside this file
mentions them. They read as symmetry-driven additions to the \texttt{inst} pair rather
than lemmas the projection development needs.}
\end{proof}

\begin{definition}[The head constant of a spine, and the type former a motive ranges over]
\label{def:headconst}
\lean{Lean4Lean.VExpr.headConst?,Lean4Lean.VExpr.motiveFormer?,Lean4Lean.VExpr.headConst?_eq_some}
\leanok
\contrib{iota} \srcloc{Lean4Lean/Theory/VExpr.lean}{1279}
\uses{def:app-spine,def:pi-telescope}
\texttt{headConst? e} is the head constant of $e$'s application spine, if the head is a
constant at all, and \texttt{headConst?\_eq\_some} reflects that option into an
existential over the level arguments. \texttt{motiveFormer? A} is the head constant of
the \emph{last} $\Pi$-binder of $A$, i.e. of the $P$ in a motive type
$\forall a{::}\alpha.\,P\,a \to \mathsf{U}$. That is the trick which lets block
well-formedness say ``this recursor eliminates that type former'' without carrying a
separate field.

\thesisref{axioms.tex \S2.6.3, the motive type $\kappa = \forall a{::}\alpha.\,P\,a\to \mathsf{U}_u$}
\end{definition}

\begin{lemma}[Constant-headed spines as decidable conditions]
\label{fam:const-spine}
\lean{Lean4Lean.VExpr.const_mkApps_spine,Lean4Lean.VExpr.eq_const_mkApps_of_spine,Lean4Lean.VExpr.eq_const_mkApps_iff,Lean4Lean.VExpr.eq_const_mkApps_append_iff}
\leanok
\contrib{iota} \srcloc{Lean4Lean/Theory/VExpr.lean}{1295}
\uses{def:headconst,def:app-spine,def:mkapps}
A constant application $(\texttt{const}\;c\;us).\texttt{mkApps}\;\mathit{args}$ has spine
head $\texttt{const}\;c\;us$ and arguments $\mathit{args}$, and conversely an expression
with those spine projections is that application. The two biconditionals package this:
$e$ is $c$ applied to exactly $\mathit{args}$ iff $\texttt{e.headConst?} = \texttt{some}\;c$
and $\texttt{e.getAppArgs} = \mathit{args}$; and $e$ is $c$ applied to $\mathit{pre}$
followed by further arguments satisfying $P$ iff the head constant is $c$, $\mathit{pre}$
is a prefix of the arguments, and $P$ holds of
$\texttt{e.getAppArgs.drop}\;|\mathit{pre}|$.

\thesisref{no thesis counterpart (a decidability-oriented reformulation)}
\end{lemma}
\begin{proof}
\uses{def:app-spine}
\leanok
Both directions from \texttt{mkApps\_getAppFn\_getAppArgs} together with
\texttt{headConst?\_eq\_some}. The last biconditional is what makes the
\texttt{ValidIndApp} and \texttt{CtorResult} clauses of block well-formedness decidable
(\ref{def:ind-valid-ind-app}, \ref{def:ind-ctor-result}).
\reviewnote{\texttt{const\_mkApps\_spine} bundles two independent facts into one
conjunction and is only ever consumed as \texttt{.1}/\texttt{.2}; two lemmas would read
better. Neither it nor \texttt{eq\_const\_mkApps\_of\_spine} is used outside
\texttt{VExpr.lean}.}
\end{proof}

\section{Environments}

\begin{definition}[Constant signatures and definitional-equality axioms]
\label{struct:vconstant}
\lean{Lean4Lean.VConstant,Lean4Lean.VDefEq}
\leanok
\uses{ind:vexpr}
\srcloc{Lean4Lean/Theory/VEnv.lean}{6}
\texttt{VConstant} is a universe-parameter count together with a type;
\texttt{VDefEq} is a universe count together with a left-hand side, a right-hand side
and their common type.

\thesisref{axioms.tex \S2.5, $\tau_{\bar u}(c)$ and the $\delta$ rule $c_{\bar\ell}\equiv v_{\bar\ell}(c)$}
\end{definition}

\begin{definition}[Environments \texttt{VEnv}]
\label{struct:venv}
\lean{Lean4Lean.VEnv}
\leanok
\contrib{mixed} \srcloc{Lean4Lean/Theory/VEnv.lean}{16}
\uses{struct:vconstant,inductive:pattern-core}
An environment has three components: a partial map \texttt{constants} from names to
constant signatures, a \emph{predicate} \texttt{defeqs} on \texttt{VDefEq} giving the
axiomatic definitional equalities, and (added on the \texttt{iota} branch) a dependent
predicate
$\texttt{pats} : (p : \texttt{Pattern}) \to p.\texttt{RHS} \times p.\texttt{Check} \to \texttt{Prop}$,
the registry of schematic reduction rules keyed by the shape of the redex. Making
reduction rules \texttt{Prop}-valued environment data rather than a fixed clause of the
reduction relation keeps environments extensible and monotone, and matches how master
already treats \texttt{defeqs}.

\thesisref{axioms.tex \S2.5 and \S2.6, the ambient set of constants and reduction rules}
\reviewnote{The dependent type of \texttt{pats} is heavier than that of the other two
fields and is what forces the transport in \ref{def:venv-addpat}.}
\end{definition}

\begin{definition}[The empty environment and its extensions]
\label{def:venv-basic}
\lean{Lean4Lean.VEnv.empty,Lean4Lean.VEnv.contains,Lean4Lean.VEnv.addConst,Lean4Lean.VEnv.addDefEq}
\leanok
\contrib{mixed} \srcloc{Lean4Lean/Theory/VEnv.lean}{23}
\uses{struct:venv}
\texttt{empty} has no constants, no definitional equalities and no pattern rules;
\texttt{contains} tests for a declared name; \texttt{addConst} returns \texttt{none} if
the name is already taken and otherwise extends the map; \texttt{addDefEq}
unconditionally adjoins one equation to the \texttt{defeqs} predicate. Only the
\texttt{pats \_ \_ := False} line of \texttt{empty} is \texttt{iota} work.

\thesisref{axioms.tex \S2.5, extending the environment by a declaration}
\end{definition}

\begin{definition}[Registering a schematic reduction rule]
\label{def:venv-addpat}
\lean{Lean4Lean.VEnv.addPat}
\leanok
\contrib{iota} \srcloc{Lean4Lean/Theory/VEnv.lean}{44}
\uses{struct:venv,inductive:pattern-rhs}
\texttt{env.addPat p r} adjoins the rule \texttt{r} for the pattern \texttt{p}: the new
\texttt{pats p' r'} holds when the old one did, or when
$\exists\,h : p' = p,\; h \triangleright r' = r$. The dependent equality is needed
because $r'$ and $r$ inhabit
$p'.\texttt{RHS} \times p'.\texttt{Check}$ and $p.\texttt{RHS} \times p.\texttt{Check}$,
which are different types.

\thesisref{axioms.tex \S2.6.4 (\texttt{sec:iota}): ``technically, the reduction rule is all substitution instances of this rule''}
\reviewnote{Unlike \texttt{addConst}, \texttt{addPat} cannot fail: there is no freshness
or consistency check, so two conflicting reducts for the same pattern are admissible at
this layer. That burden is pushed entirely onto block well-formedness and
\texttt{VEnv.PatWF} downstream, and the asymmetry with \texttt{addConst} deserves a
comment here. The idiom $\exists\,h : p' = p,\; h \triangleright r' = r$ is correct but
unusual; a heterogeneous equality would read more conventionally.}
\end{definition}

\begin{definition}[Environment extension order]
\label{struct:venv-le}
\lean{Lean4Lean.VEnv.LE,Lean4Lean.VEnv.LE.rfl,Lean4Lean.VEnv.LE.trans}
\leanok
\contrib{mixed} \srcloc{Lean4Lean/Theory/VEnv.lean}{47}
\uses{struct:venv}
$\texttt{env}_1 \le \texttt{env}_2$ is a three-field structure: every constant lookup,
every definitional equality and (the \texttt{iota} addition, line 50) every registered
pattern rule of $\texttt{env}_1$ is present in $\texttt{env}_2$. Reflexivity is the
triple of identities and transitivity the triple of composites, so extending the order
with the new field cost two lines. This is the monotonicity backbone on which every
environment-extension lemma rests.

\thesisref{no thesis counterpart (an internal invariant)}
\end{definition}

\section{Declarations}

\begin{definition}[Named constants and definitions]
\label{def:vconstval}
\lean{Lean4Lean.VConstVal,Lean4Lean.VDefVal,Lean4Lean.VDefVal.toDefEq}
\leanok
\uses{struct:vconstant}
\srcloc{Lean4Lean/Theory/VDecl.lean}{5}
\texttt{VConstVal} is a \texttt{VConstant} with a name; \texttt{VDefVal} adds a value.
\texttt{toDefEq} turns a definition into the $\delta$ equation
$c_{\texttt{params uvars}} \equiv v$ at the declared type.

\thesisref{axioms.tex \S2.5 (definitions, $\delta$ reduction)}
\end{definition}

\begin{definition}[One type former of an inductive block]
\label{struct:vinductivetype}
\lean{Lean4Lean.VInductiveType}
\leanok
\uses{def:vconstval}
\srcloc{Lean4Lean/Theory/VDecl.lean}{14}
A named constant together with the list of its constructors, each again a named constant
with a type.

\thesisref{axioms.tex \S2.6.1, $K = \sum_c (c : \forall b{::}\beta.\,t\;p[b])$}
\end{definition}

\begin{definition}[One recursor computation rule]
\label{struct:vrecrule}
\lean{Lean4Lean.VRecRule}
\leanok
\contrib{iota} \srcloc{Lean4Lean/Theory/VDecl.lean}{20}
\uses{ind:vexpr}
A rule records the constructor \texttt{ctor} it fires on, that constructor's parameter
count \texttt{ctorParams}, its number of non-parameter arguments \texttt{nfields}, and
the closed reduct template \texttt{rhs}; the redex key's constructor spine therefore has
$\texttt{ctorParams} + \texttt{nfields}$ arguments. It mirrors
\texttt{Lean.RecursorRule} except for \texttt{ctorParams}, which the kernel's structure
does not have.

\thesisref{axioms.tex \S2.6.4 (\texttt{sec:iota}), the rule $\mathrm{rec}_P\,C\,e\,p[b]\,(c\,b) \rightsquigarrow e_c\,b\,v$}
\reviewnote{The docstring justifies the extra field carefully: for a nested inductive the
constructor a rule fires on need not have the recursor's parameter count
(\texttt{Tree.rec\_1} fires on \texttt{List.cons}), and it is honest that the current
well-formedness predicate admits only direct blocks, where the two coincide. The field
is therefore redundant as data but necessary as specification, and the code says so.}
\end{definition}

\begin{definition}[A recursor declaration]
\label{struct:vrecursor}
\lean{Lean4Lean.VRecursor}
\leanok
\contrib{iota} \srcloc{Lean4Lean/Theory/VDecl.lean}{35}
\uses{struct:vrecrule,def:vconstval}
A named constant extended with the block it belongs to (\texttt{all}), the telescope
segmentation \texttt{numParams}, \texttt{numMotives}, \texttt{numMinors},
\texttt{numIndices}, the K-like-reduction flag \texttt{k}, and its list of computation
rules. It mirrors \texttt{Lean.RecursorVal} field for field, which is what makes the
kernel-side correspondence a field-by-field match.

\thesisref{axioms.tex \S2.6.3 (the recursor), the segmentation $\forall C{:}\kappa.\,\forall e{::}\varepsilon.\,\forall a{::}\alpha.\,\forall z{:}P\,a$}
\reviewnote{\texttt{k} is recorded but, as its own docstring admits, unused by the
theory: the thesis's second $\iota$ rule (K-like reduction on a subsingleton eliminator)
is outside the model. The honesty is in the docstring, but the structure still advertises
coverage the development does not have.}
\end{definition}

\begin{definition}[Positions in the recursor telescope]
\label{def:getmajoridx}
\lean{Lean4Lean.VRecursor.getMajorIdx,Lean4Lean.VRecursor.getFirstIndexIdx}
\leanok
\contrib{iota} \srcloc{Lean4Lean/Theory/VDecl.lean}{46}
\uses{struct:vrecursor}
\texttt{getMajorIdx} is the argument index of the major premise,
$\texttt{numParams}+\texttt{numMotives}+\texttt{numMinors}+\texttt{numIndices}$, and
\texttt{getFirstIndexIdx} that of the first index; both mirror the homonymous fields of
\texttt{Lean.RecursorVal}.

\thesisref{axioms.tex \S2.6.3, the position of the major premise $z$ in the recursor telescope}
\reviewnote{\texttt{VRecursor.getFirstIndexIdx} has no consumer anywhere. The
\texttt{getFirstIndexIdx} occurrences elsewhere in the repository are all Lean's own
\texttt{RecursorVal} field on a kernel \texttt{rval}
(\srcloc{Lean4Lean/Inductive/Reduce.lean}{84} and
\srcloc{Lean4Lean/Verify/TypeChecker/WHNF.lean}{29}), plus one mention in a comment in
\texttt{Theory/Typing/Pattern.lean}. And \texttt{VEnv.addRecRule}
(\srcloc{Lean4Lean/Theory/Inductive.lean}{260}) spells the major index
$\texttt{numParams}+\texttt{numMotives}+\texttt{numMinors}+\texttt{numIndices}$ out by
hand instead of calling \texttt{getMajorIdx}, although its docstring says ``major at
\texttt{getMajorIdx}''.}
\end{definition}

\begin{definition}[An inductive block]
\label{struct:vinductdecl}
\lean{Lean4Lean.VInductDecl}
\leanok
\contrib{mixed} \srcloc{Lean4Lean/Theory/VDecl.lean}{54}
\uses{struct:vinductivetype,struct:vrecursor}
A mutual inductive declaration: the number of universe parameters, the number of type
parameters, the list of type formers with their constructors, and (the \texttt{iota}
addition, line 58) the list of recursors with their computation rules. Storing
recursors as \emph{data} rather than deriving them is the central design choice of the
\texttt{iota} contribution: it mirrors the kernel, which stores them, at the cost of
having to \emph{specify} rather than construct them, which is what block well-formedness
(\ref{structure:ind-wf}) then does.

\thesisref{axioms.tex \S2.6.1, $\mu t{:}F.\;K$, together with \S2.6.3's $\mathrm{rec}_P$}
\end{definition}

\begin{definition}[Declaration kinds]
\label{ind:vdecl}
\lean{Lean4Lean.VDecl}
\leanok
\uses{def:vconstval,struct:vinductdecl}
\srcloc{Lean4Lean/Theory/VDecl.lean}{60}
The ways an environment may be extended: \texttt{axiom}, \texttt{def},
\texttt{opaque}, \texttt{example}, the quotient block, an inductive block, or a mutual
block of definitions.
\end{definition}

\section{Translating real Lean syntax}

\begin{definition}[Translating \texttt{Lean.Expr} into \texttt{VExpr}]
\label{def:meta-ofexpr}
\lean{Lean4Lean.Meta.expandProj,Lean4Lean.Meta.expandExpr,Lean4Lean.Meta.ofExpr}
\leanok
\uses{ind:vexpr,def:vlevel-oflevel}
\srcloc{Lean4Lean/Theory/Meta.lean}{12}
\texttt{expandProj} rewrites a structure projection into the corresponding
\texttt{casesOn} application, \texttt{expandExpr} does so everywhere and instantiates
metavariables, and \texttt{ofExpr} maps the residual syntax into \texttt{VExpr},
translating \texttt{let} by substitution into the body, literals through
\texttt{toConstructor}, and free variables through a supplied index map or their local
definition. Projections and metavariables that survive are rejected with an error.

\reviewnote{This is the context for the \texttt{trproj} work: the model's syntax has no
projection node, and this translator removes projections by \texttt{casesOn} expansion,
so projections must be modelled as a reduction relation elsewhere. Defensible, but it
means \texttt{VExpr} and \texttt{Lean.Expr} diverge on exactly the construct being
verified.}
\end{definition}

\begin{definition}[The \texttt{vexpr}/\texttt{vconst}/\texttt{vdefeq} term macros]
\label{def:meta-macros}
\lean{Lean4Lean.Meta.toVExprCore,Lean4Lean.Meta.toVExprWrapper,Lean4Lean.Meta.elabForVExpr}
\leanok
\uses{def:meta-ofexpr}
\srcloc{Lean4Lean/Theory/Meta.lean}{59}
Elaboration-time macros that let a \texttt{VExpr}, a \texttt{VConstant} or a
\texttt{VDefEq} be written as ordinary Lean syntax and compiled to its model
representation, the auto-bound universe parameters being counted into \texttt{uvars}.
They are used by \texttt{Theory/Quot.lean} and by the primitive-constant tables of the
\texttt{Verify} layer.
\end{definition}

\section{The quotient block}

\begin{definition}[The quotient constants and their computation rule]
\label{def:quot-consts}
\lean{Lean4Lean.eqConst,Lean4Lean.quotConst,Lean4Lean.quotMkConst,Lean4Lean.quotLiftConst,Lean4Lean.quotIndConst,Lean4Lean.quotDefEq}
\leanok
\uses{def:meta-macros,struct:vconstant}
\srcloc{Lean4Lean/Theory/Quot.lean}{6}
The signatures of \texttt{Eq}, \texttt{Quot}, \texttt{Quot.mk}, \texttt{Quot.lift} and
\texttt{Quot.ind} are obtained by elaborating the real Lean constants through
\texttt{type\_of\%} and the \texttt{vconst} macro, and \texttt{quotDefEq} is the equation
$\mathrm{lift}_R\,\beta\,f\,h\,(\mathrm{mk}_R\,a) \equiv f\,a$. \texttt{Quot.sound} is
absent here, correctly: in Lean 4 it is an ordinary \texttt{axiom} and not one of the
four quotient primitives the kernel's quotient module installs, so in the model it
arrives through the generic \texttt{VDecl.axiom} path rather than through
\texttt{addQuot}. The kernel-side \texttt{Environment.addQuot}
(\srcloc{Lean4Lean/Quot.lean}{39}) checks the same four names.

\thesisref{axioms.tex \S2.7.1 (quotient types)}
\reviewnote{Taking the types from real Lean rather than writing them out is economical,
but it makes the model's fidelity depend on the ambient Lean version.}
\end{definition}

\begin{definition}[Adding the quotient block to an environment]
\label{def:venv-addquot}
\lean{Lean4Lean.VEnv.QuotReady,Lean4Lean.VEnv.addQuot}
\leanok
\uses{def:quot-consts,def:venv-basic}
\srcloc{Lean4Lean/Theory/Quot.lean}{13}
\texttt{QuotReady env} requires \texttt{Eq} to be already declared with the expected
signature; \texttt{addQuot} then adds the four quotient constants in order and finally
the \texttt{Quot.lift}/\texttt{Quot.mk} definitional equality, failing if any name is
taken.

\thesisref{axioms.tex \S2.7.1}
\reviewnote{The quotient computation rule is registered as a \texttt{VDefEq}, not through
the \texttt{pats} registry of \ref{struct:venv}, so after the \texttt{iota} work the
model carries two mechanisms for computation rules. This is a follow-up rather than a
defect, but nothing in the code says so: \texttt{Theory/Quot.lean} has no docstrings at
all, and neither \texttt{VEnv.addPat} nor \texttt{Theory/Typing/Pattern.lean} mentions
the quotient rule.}
\end{definition}

\section{Level equivalence is coNP-hard}

This file has no counterpart in the thesis. It is a self-contained reduction from CNF
satisfiability to \texttt{VLevel} equivalence, and hence a lower bound on the problem
that the kernel's \texttt{Level.isEquiv} and \texttt{Level.geq} solve. The module
docstring also sketches the matching NP upper bound for the complement, which is not
formalised.

\begin{definition}[CNF formulas and the dual-rail encoding]
\label{def:levelsat-cnf}
\lean{Lean4Lean.LevelSat.Lit,Lean4Lean.LevelSat.Clause,Lean4Lean.LevelSat.CNF,Lean4Lean.LevelSat.ClauseSat,Lean4Lean.LevelSat.Sat,Lean4Lean.LevelSat.yes,Lean4Lean.LevelSat.no,Lean4Lean.LevelSat.rail}
\leanok
\srcloc{Lean4Lean/Theory/LevelSat.lean}{70}
Literals are pairs in $\mathrm{Fin}\,n \times \mathrm{Bool}$, clauses are lists of
literals, a CNF is a list of clauses, and \texttt{Sat} is the usual existential over
Boolean assignments. Each variable $i$ gets two level parameters, the rails $i^{+}$ at
index $2i$ and $i^{-}$ at index $2i+1$, a rail counting as ``true'' when its value is
nonzero. Indexing variables by $\mathrm{Fin}\,n$ puts ``the formula mentions only
variables the levels provide rails for'' into the type, so the reduction carries no side
condition.
\end{definition}

\begin{definition}[The two levels of the reduction]
\label{def:levelsat-levels}
\lean{Lean4Lean.LevelSat.two,Lean4Lean.LevelSat.clauseLvl,Lean4Lean.LevelSat.gates,Lean4Lean.LevelSat.var,Lean4Lean.LevelSat.small,Lean4Lean.LevelSat.big}
\leanok
\uses{def:levelsat-cnf,ind:vlevel}
\srcloc{Lean4Lean/Theory/LevelSat.lean}{104}
$\mathrm{clauseLvl}\,C$ is the $\max$ of the rails of a clause; $\mathrm{gates}\,\varphi$
gates the constant $2$ behind every clause by an iterated $\mathrm{imax}$ fold;
$\mathrm{var}\,i = \max(i^{+},\,\mathrm{imax}(\mathrm{imax}(2,i^{+}),\,i^{-}))$ detects a
variable whose two rails are both nonzero; $\mathrm{small}\,n$ is the $\max$ of all
$\mathrm{var}\,i$ for $i < n$; and $\mathrm{big}\,\varphi = \max(\mathrm{small}\,n,\,\mathrm{gates}\,\varphi)$.
Since every level denotes a monotone function, negation cannot be computed; it enters as
a failure of domination, the constant $2$ having to exceed everything the right-hand
side offers.
\end{definition}

\begin{lemma}[Evaluation lemmas for the gadget levels]
\label{fam:levelsat-eval}
\lean{Lean4Lean.LevelSat.eval_var_eq,Lean4Lean.LevelSat.yes_le_eval_var,Lean4Lean.LevelSat.no_le_eval_var,Lean4Lean.LevelSat.two_le_eval_var,Lean4Lean.LevelSat.eval_var_le_one,Lean4Lean.LevelSat.le_eval_small,Lean4Lean.LevelSat.two_le_eval_small,Lean4Lean.LevelSat.eval_small_le_one,Lean4Lean.LevelSat.eval_clauseLvl_ne_zero,Lean4Lean.LevelSat.eval_foldl_imax_le,Lean4Lean.LevelSat.eval_foldl_imax_le_of_zero,Lean4Lean.LevelSat.le_eval_foldl_imax}
\leanok
\uses{def:levelsat-levels,def:vlevel-eval}
\srcloc{Lean4Lean/Theory/LevelSat.lean}{141}
Pointwise evaluation of the gadgets. $\mathrm{var}\,i$ is the larger of the two rails
unless both are nonzero, in which case it is at least $2$ (\texttt{eval\_var\_eq} is the
equation, the four inequalities its consequences). For $\mathrm{small}\,n$ there are
three separate one-directional lemmas, not a biconditional: it dominates every parameter
below $2n$ (\texttt{le\_eval\_small}), it is at least $2$ as soon as both rails of some
$i<n$ are nonzero (\texttt{two\_le\_eval\_small}), and it is at most $1$ given that every
$i<n$ has a zero rail and every parameter below $2n$ is at most $1$
(\texttt{eval\_small\_le\_one}). The $\mathrm{imax}$ fold is bounded by the gate values
unless every gate is nonzero, in which case the gated constant survives.
\end{lemma}
\begin{proof}
\uses{def:vlevel-eval}
\leanok
Each lemma restates \texttt{VLevel.eval} in terms of \texttt{max} and \texttt{ite} and
closes with \texttt{split} and \texttt{omega}; the \texttt{small} lemmas are inductions
on $n$ and the fold lemmas inductions on the list of gates.
\end{proof}

\begin{definition}[The rail encoding of an assignment]
\label{def:levelsat-raillist}
\lean{Lean4Lean.LevelSat.railVal,Lean4Lean.LevelSat.railList,Lean4Lean.LevelSat.ofFin,Lean4Lean.LevelSat.getD_railList,Lean4Lean.LevelSat.getD_railList_rail}
\leanok
\uses{def:levelsat-cnf}
\srcloc{Lean4Lean/Theory/LevelSat.lean}{270}
$\mathrm{railVal}\,a\,j$ is $1$ when rail $j$ agrees with the assignment $a$ and $0$
otherwise, and $\mathrm{railList}\,n\,a$ is the list of the first $2n$ rail values: the
point at which $\mathrm{big}$ and $\mathrm{small}$ differ when $a$ satisfies the formula.
\texttt{ofFin} reads a $\mathrm{Fin}\,n$-indexed assignment at an arbitrary index.
\end{definition}

\begin{theorem}[Level equivalence is coNP-hard]
\label{thm:levelsat-equiv-iff-unsat}
\lean{Lean4Lean.LevelSat.equiv_iff_unsat}
\leanok
\uses{def:levelsat-levels,def:levelsat-cnf,def:vlevel-le-equiv}
\srcloc{Lean4Lean/Theory/LevelSat.lean}{331}
For every CNF formula $\varphi$ over $n$ variables,
$\mathrm{big}\,\varphi \approx \mathrm{small}\,n$ if and only if $\varphi$ is
unsatisfiable. Since the two levels are computed from $\varphi$ in linear size
(\ref{fam:levelsat-size}), deciding \texttt{VLevel.Equiv}, the problem
\texttt{Level.isEquiv} solves, is coNP-hard.
\end{theorem}
\begin{proof}
\uses{fam:levelsat-eval,def:levelsat-raillist}
\leanok
Forward: a satisfying assignment gives the point $\mathrm{railList}\,n\,(\mathrm{ofFin}\,a)$
at which $\mathrm{small}$ stays at most $1$, the rails being $0/1$ and consistent, while
every clause gate is nonzero, so the gated constant $2$ pokes above it. Backward: at any
point where the two levels differ, the gate fold must exceed $\mathrm{small}$, which
forces every clause gate nonzero and every variable's rails consistent; reading the
assignment off the $\mathrm{yes}$ rails yields a satisfying assignment.
\end{proof}

\begin{lemma}[The same hardness for the order test]
\label{thm:levelsat-le-iff-unsat}
\lean{Lean4Lean.LevelSat.le_iff_unsat,Lean4Lean.LevelSat.small_le_big}
\leanok
\uses{def:levelsat-levels,def:vlevel-le-equiv}
\srcloc{Lean4Lean/Theory/LevelSat.lean}{382}
$\mathrm{big}\,\varphi \le \mathrm{small}\,n$ if and only if $\varphi$ is unsatisfiable,
so \texttt{Level.geq} is coNP-hard as well.
\end{lemma}
\begin{proof}
\uses{thm:levelsat-equiv-iff-unsat}
\leanok
Rewrite by \ref{thm:levelsat-equiv-iff-unsat} and antisymmetry
(\texttt{VLevel.le\_antisymm\_iff}); the missing half is
$\mathrm{small}\,n \le \mathrm{big}\,\varphi$, which holds structurally by
\texttt{Nat.le\_max\_left}.
\end{proof}

\begin{lemma}[The reduction is linear]
\label{fam:levelsat-size}
\lean{Lean4Lean.LevelSat.size,Lean4Lean.LevelSat.cnfSize,Lean4Lean.LevelSat.size_clauseLvl,Lean4Lean.LevelSat.size_gates,Lean4Lean.LevelSat.size_var,Lean4Lean.LevelSat.size_small,Lean4Lean.LevelSat.size_big}
\leanok
\uses{def:levelsat-levels,def:levelsat-cnf}
\srcloc{Lean4Lean/Theory/LevelSat.lean}{392}
Counting nodes, $\mathrm{size}(\mathrm{big}\,\varphi) = 10n + 2\,\mathrm{cnfSize}(\varphi) + 5$,
where $\mathrm{cnfSize}$ charges one for each literal and one for each clause. So both
levels are linear in the formula, which is what makes
\ref{thm:levelsat-equiv-iff-unsat} a hardness result.
\end{lemma}
\begin{proof}
\uses{def:levelsat-levels}
\leanok
Inductions on the clause, the clause list and $n$, each closing with \texttt{omega};
$\mathrm{size}(\mathrm{var}\,i) = 9$ and $\mathrm{size}(2) = 3$ are \texttt{rfl}. The
docstring is explicit that only the size half of the linear-time claim is formalised.
\end{proof}

\begin{example}[Two worked instances]
\label{ex:levelsat-instances}
\lean{Lean4Lean.LevelSat.selfContradiction,Lean4Lean.LevelSat.not_sat_selfContradiction}
\leanok
\uses{thm:levelsat-equiv-iff-unsat}
\srcloc{Lean4Lean/Theory/LevelSat.lean}{431}
$x_0 \wedge \neg x_0$ is unsatisfiable, so the two levels it produces are equivalent;
this is exactly the case where a naive reduction reading each parameter directly as a SAT
variable would wrongly report a difference, since the rails $0^{+}$ and $0^{-}$ may both
be nonzero and $\mathrm{small}$ is what rules that point out. The formula $x_0$ alone is
satisfiable, and the two levels differ at the rail encoding $[1,0]$, where
$\mathrm{big}$ evaluates to $2$ and $\mathrm{small}$ to $1$.
\end{example}

\section{Contribution summary}

The \texttt{iota} branch contributed 254 lines here, and they are specification design
rather than proof. Three pieces matter. First, the \texttt{pats} field of
\ref{struct:venv} with \ref{def:venv-addpat} and the third clause of
\ref{struct:venv-le}: reduction rules become \texttt{Prop}-valued, pattern-indexed
environment data instead of a fixed clause of the reduction relation, which keeps
environments extensible and monotone; in the monotonicity backbone it costs one new
field (\srcloc{Lean4Lean/Theory/VEnv.lean}{50}) and two rewritten proof terms, the
pairs of \texttt{LE.rfl} and \texttt{LE.trans} becoming triples.
Second, \ref{struct:vrecrule}, \ref{struct:vrecursor}, \ref{def:getmajoridx} and the
\texttt{recs} field of \ref{struct:vinductdecl}: recursors and their computation rules
become data in the model, mirroring \texttt{Lean.RecursorVal} and
\texttt{Lean.RecursorRule} field for field, with the deliberate addition of
\texttt{ctorParams}. Third, the syntactic toolkit at the end of \texttt{VExpr.lean}
(\ref{def:mkapps} to \ref{fam:const-spine}): total telescope and spine readers, plus
biconditionals trading existentials over expression lists for prefix and length tests,
which is what lets the shape predicates of
chapter \ref{chap:inductive} be decidable and therefore testable against the real kernel.
The one load-bearing lemma of that block is the disjointness in \ref{def:ctorheaded}.

The \texttt{trproj} branch contributed 195 lines, all in \texttt{VExpr.lean}, and they
are supporting infrastructure for the projection development of chapter
\ref{chap:proj}. The organising idea is \ref{thm:lift-eq-subst}: since weakening and
instantiation are both substitutions, a builder lemma need only be proved once against
\texttt{subst}. That idea is written down as a section docstring and pays off
immediately in \ref{fam:mkapps-stability}, where one induction yields three corollaries.
What is actually consumed downstream splits in two. \texttt{Theory/Proj.lean} takes
\ref{thm:lift-eq-subst}, all five lemmas of \ref{fam:subst-liftn-algebra},
\ref{thm:subst-instn}, \ref{thm:liftn-subst-liftn}, the \texttt{subst}, \texttt{inst} and
\texttt{instL} halves of \ref{fam:mkapps-stability}, \texttt{getElem\_bvarsDesc}, and
\texttt{CtorHeaded.inst} and \texttt{piArity\_inst\_of\_ctorHeaded} from
\ref{fam:ctorheaded-stability}. The \emph{verification} layer, not
\texttt{Theory/Proj.lean}, is what consumes \ref{fam:pibinders-stability}
(\srcloc{Lean4Lean/Verify/Typing/Lemmas.lean}{652},
\srcloc{Lean4Lean/Verify/Typing/Lemmas.lean}{1598},
\srcloc{Lean4Lean/Verify/Environment/Lemmas.lean}{1051}),
\texttt{piBinders\_inst\_of\_ctorHeaded}, \texttt{mkApps\_lift'}, and the
$\lambda$-telescope readers of \ref{def:lam-telescope}.
\ref{thm:lift-instn-hi} is consumed only through its $m=0$ case, which
\texttt{Theory/Typing/HeadReduction.lean} already used before the refactor.

What remains open at this layer. The K-like reduction flag of \ref{struct:vrecursor} is
recorded but not modelled, so the thesis's second $\iota$ rule is out of scope. The
quotient computation rule (\ref{def:venv-addquot}) still lives in \texttt{defeqs} rather
than in \texttt{pats}, leaving two mechanisms for computation rules. \texttt{ctorParams}
is redundant as data until nested blocks are admitted. And seven contributed
declarations in this chapter are dead. Five are mentioned nowhere outside their own
statement: \texttt{Lift.liftVar\_consN\_lt}, \texttt{Lift.liftVar\_consN\_succ},
\texttt{VExpr.CtorHeaded.forallE}, \texttt{VExpr.CtorHeaded.instL} and
\texttt{VRecursor.getFirstIndexIdx}. Two more, \texttt{Lift.consN\_fixes} and
\texttt{VExpr.getAppFn\_instL\_const}, survive only by feeding one of those five. The
review notes below list them with line numbers.

\section{Review notes}

\begin{itemize}
\item \textbf{Medium.} \srcloc{Lean4Lean/Theory/VExpr.lean}{919}: dead upstream code
created by the \texttt{trproj} refactor. Master proved \texttt{lift'\_inst\_hi} at
\texttt{master:VExpr.lean:903} with the one-line
\texttt{simp [subst\_lift', lift'\_subst, lift\_r\_one, inst\_eq]}; the branch moved the
lemma to line 1018 and re-derived it from the new \texttt{lift'\_instN\_hi}.
\texttt{lift\_r\_one} now has no consumer anywhere, and with it the whole chain that
existed to feed it: \texttt{Subst.lift\_r\_comm} (911), \texttt{Subst.trunc} (908),
\texttt{Subst.Depth.one} (885), \texttt{Subst.Depth.id} (858) and \texttt{Subst.Depth}
(796) are now mentioned only inside that block. Six declarations. Either keeping
master's proof or deleting the orphaned block would be consistent; leaving both
maximises merge-conflict surface on a file otherwise shared verbatim with upstream
(\ref{fam:subst-trunc-orphans}, \ref{thm:lift-instn-hi}).
\item \textbf{Medium.} \srcloc{Lean4Lean/Theory/VExpr.lean}{1232}:
\texttt{VExpr.CtorHeaded.forallE} is unused anywhere in the repository, is proved by
\texttt{:= h} because it is definitionally trivial, and its docstring states a different
proposition than the lemma. The docstring text is the motivation for
\texttt{piArity\_inst\_of\_ctorHeaded}, not for this statement
(\ref{thm:ctorheaded-forallE}).
\item \textbf{Low.} \srcloc{Lean4Lean/Theory/VExpr.lean}{1046}: stale section docstring.
It promises telescope and spine readers ``with their decidability'', but the decidability
instances were relocated to \texttt{Lean4Lean/Tests/ShapeDecide.lean}, so the section
contains no decidability at all.
\item \textbf{Low.} \srcloc{Lean4Lean/Theory/VExpr.lean}{753}:
\texttt{Lift.liftVar\_consN\_lt} and \texttt{Lift.liftVar\_consN\_succ} (line 756) have no
consumer, which makes \texttt{Lift.consN\_fixes} (line 745) dead in effect as well. Of
the four-lemma block at lines 744 to 758 only \texttt{consN\_cons} is used
(\ref{fam:lift-consn}).
\item \textbf{Low.} \srcloc{Lean4Lean/Theory/VExpr.lean}{1256}:
\texttt{VExpr.CtorHeaded.instL} and its helper \texttt{getAppFn\_instL\_const} (line 1248)
are dead; nothing outside \texttt{VExpr.lean} mentions them. They look like
symmetry-driven additions to the \texttt{inst} pair rather than lemmas the projection
development needs (\ref{fam:ctorheaded-stability}).
\item \textbf{Low.} \srcloc{Lean4Lean/Theory/VExpr.lean}{1009}: redundancy left
unresolved, and inconsistently. \texttt{subst\_instN} subsumes master's
\texttt{subst\_inst} (line 952) at $n = 0$, and \texttt{liftN\_subst\_liftN} (line 1032)
subsumes master's \texttt{lift\_subst\_lift} (line 947) at $i = 0$ -- its own docstring
says so -- yet both master lemmas keep independent hand proofs, whereas
\texttt{lift'\_inst\_hi} in the same block \emph{was} re-derived. Three analogous pairs,
three different treatments. Additionally \texttt{subst\_instN}'s docstring says ``under
\texttt{m} binders'' while the statement binds $n$
(\ref{thm:subst-instn}, \ref{thm:liftn-subst-liftn}).
\item \textbf{Low.} \srcloc{Lean4Lean/Theory/VExpr.lean}{1133}: the whole
$\lambda$-telescope block -- \texttt{lamBinders}, \texttt{lamBinders\_length},
\texttt{foldr\_lam\_lamBinders} -- has exactly one consumer, a single proof at
\srcloc{Lean4Lean/Verify/Environment/Lemmas.lean}{1128}, and
\texttt{lamBinders\_length} is reached only through \texttt{simp} rather than by name.
That is thin, but it is not dead code: the claim that it is should not be made
(\ref{def:lam-telescope}).
\item \textbf{Low.} \srcloc{Lean4Lean/Theory/VExpr.lean}{1295}:
\texttt{const\_mkApps\_spine} bundles two independent facts into one conjunction and is
only ever consumed as \texttt{.1}/\texttt{.2}; together with
\texttt{eq\_const\_mkApps\_of\_spine} it has no consumer outside \texttt{VExpr.lean}
(\ref{fam:const-spine}).
\item \textbf{Low.} \srcloc{Lean4Lean/Theory/VExpr.lean}{1200}: the docstring
``Head-constructor tests, with their reflections into existentials'' is attached to
\texttt{isSort} alone while describing all six declarations; and the whole block exists
only for the decidability instances in the test file (\ref{def:head-tests}).
\item \textbf{Low.} \srcloc{Lean4Lean/Theory/VExpr.lean}{108}: \texttt{decClosedN} stayed
in the theory file while the sibling head-shape instances moved to
\texttt{Lean4Lean/Tests/ShapeDecide.lean} (commit \texttt{ba118fd}). That file's module
docstring explains the general policy but not this exception, which is the part a reader
needs: \texttt{decClosedN} is the one such instance a \emph{definition} depends on, via
the \texttt{if h : ru.rhs.Closed} of \texttt{VEnv.addRecRule}
(\srcloc{Lean4Lean/Theory/Inductive.lean}{258}) (\ref{inst:dec-closedn}).
\item \textbf{Low.} \srcloc{Lean4Lean/Theory/VDecl.lean}{51}:
\texttt{VRecursor.getFirstIndexIdx} is defined and documented but has no consumer
anywhere; every other \texttt{getFirstIndexIdx} in the repository is Lean's own
\texttt{RecursorVal} field on a kernel \texttt{rval}
(\srcloc{Lean4Lean/Inductive/Reduce.lean}{84},
\srcloc{Lean4Lean/Verify/TypeChecker/WHNF.lean}{29} and following), plus one comment in
\texttt{Theory/Typing/Pattern.lean}. Relatedly,
\texttt{VEnv.addRecRule} (\srcloc{Lean4Lean/Theory/Inductive.lean}{260}) spells
\texttt{numParams + numMotives + numMinors + numIndices} out by hand instead of calling
\texttt{VRecursor.getMajorIdx}, though its docstring says ``major at
\texttt{getMajorIdx}'' (\ref{def:getmajoridx}).
\item \textbf{Low.} \srcloc{Lean4Lean/Theory/VDecl.lean}{41}: \texttt{VRecursor.k} records
the K-like-reduction flag but, as its own docstring admits, is unused by the theory; the
thesis's second $\iota$ rule is outside the model. Honest, but the structure advertises
coverage the development does not have, and a reader auditing ``does this model Lean's
$\iota$?'' has to read the docstring to find out (\ref{struct:vrecursor}).
\item \textbf{Low.} \srcloc{Lean4Lean/Theory/VEnv.lean}{44}: \texttt{VEnv.addPat} cannot
fail. Unlike \texttt{addConst} it performs no freshness or consistency check, so
conflicting reducts for the same pattern are admissible at this layer; the invariant is
enforced only downstream. The asymmetry with \texttt{addConst} deserves a comment here
(\ref{def:venv-addpat}).
\item \textbf{Low.} \srcloc{Lean4Lean/Theory/Quot.lean}{11}: after the \texttt{iota} work
the model has two mechanisms for computation rules, since \texttt{quotDefEq} is still a
\texttt{VDefEq} rather than a \texttt{pats} entry. The duplication is nowhere
acknowledged in the code -- \texttt{Theory/Quot.lean} carries no docstrings, and neither
\texttt{VEnv.addPat} nor \texttt{Theory/Typing/Pattern.lean} mentions the quotient rule
(\ref{def:venv-addquot}).
\item \textbf{Low.} \srcloc{Lean4Lean/Theory/Quot.lean}{6}: the quotient signatures are
obtained by elaborating the real Lean constants, so the model's fidelity at this point
depends on the ambient Lean version rather than on anything written in the repository
(\ref{def:quot-consts}).
\item \textbf{Low (context, not a defect).} \srcloc{Lean4Lean/Theory/Meta.lean}{12}:
\texttt{ofExpr} rejects \texttt{.proj} (line 46) and \texttt{expandProj} eliminates
projections through \texttt{casesOn} beforehand. \texttt{VExpr} and \texttt{Lean.Expr}
therefore diverge on exactly the construct the \texttt{trproj} branch verifies, which is
why projections are modelled as a reduction relation in chapter \ref{chap:proj} instead
of as syntax (\ref{def:meta-ofexpr}).
\end{itemize}
